% =========================================================
% CAPITOLO 3 - METODOLOGIA
% =========================================================

\chapter{Metodologia}
\label{chap:metodologia}

La struttura del \textit{framework} metodologico per il \textit{forecasting} predittivo della
\textit{non-compliance} in sistemi distribuiti a microservizi è ibrida; combina infatti due
formalismi distinti, con ruoli complementari e non intercambiabili. I Grafi Temporali
Attribuiti (\gls{ATG}) modellano l'evoluzione dinamica delle metriche osservabili a
\textit{runtime} e forniscono la struttura dati ideale per l'applicazione dei modelli di
previsione quantitativa sulle serie temporali. Tuttavia, essi non permettono di esprimere
nativamente i vincoli logici che legano molteplici componenti a una specifica proprietà
certificata. Gli ipergrafi, invece, codificano tali vincoli logici con assoluta precisione
strutturale, ma sono per natura privi di capacità predittive. A questi due formalismi si
aggiunge un terzo elemento, la funzione di \textit{mapping} semantico, che funge da ponte
traducendo le previsioni quantitative locali in valutazioni globali di conformità sulle
proprietà certificate. 
L'architettura ibrida nasce dalla necessità di superare i limiti che un approccio standard
incontrerebbe nel modellare sistemi complessi caratterizzati da proprietà certificate multiple e
componenti condivisi tra percorsi distinti. Come illustrato nella Figura~\ref{fig:framework_overview},
il modello si articola su due dimensioni.

La prima dimensione riguarda la \textbf{rappresentazione del sistema}, organizzata su tre
\textit{layer} concettuali:

\begin{itemize}
  \item \textbf{Layer 1 (Ipergrafo di certificazione):} è una struttura statica che viene
        definita a \textit{design-time} e che permette di codificare ogni proprietà certificata
        in un iperarco con semantica AND atomica. Attraverso la funzione \textit{Shared} permette
        di identificare topologicamente i punti di interferenza tra proprietà distinte, senza
        necessità di elaborare dati storici;
  \item \textbf{Layer 2 (\textit{Attributed Temporal Graph}):} è una struttura dinamica che
        descrive l'evoluzione del sistema a \textit{runtime}, attraverso serie temporali di
        metriche osservabili. Queste serie di metriche costituiscono lo stato interno dei nodi
        e la qualità delle interazioni sugli archi. La struttura costituisce il substrato
        computazionale su cui operano tutti i modelli di previsione;
  \item \textbf{Layer 3 (\textit{Mapping} semantico):} funge da collegamento bidirezionale tra
        i due livelli precedenti. Traduce ogni \textit{compliance set} nell'insieme degli archi
        rilevanti per quella proprietà e associa le previsioni locali alla corretta funzione di
        aggregazione.
\end{itemize}

La seconda dimensione riguarda la \textbf{\textit{pipeline} di \textit{forecasting}},
articolata in quattro fasi operative atte a trasformare i dati di telemetria grezza in giudizi
di conformità.

\begin{itemize}
  \item \textbf{Fase I (Previsione locale):} la configurazione topologica dell'ipergrafo orienta
        in modo deterministico la selezione delle metriche rilevanti, mentre i modelli predittivi
        generano le traiettorie di previsione sull'orizzonte temporale configurato, elaborando le 
        serie storiche delle metriche selezionate per identificarne trend, stagionalità e dinamiche evolutive;
  \item \textbf{Fase II (Analisi causale):} una sequenza di test statistici applicati sulle coppie
        di metriche semanticamente rilevanti definisce il grafo causale, includendo le dipendenze
        \textit{cross-property} tra \textit{compliance sets} distinti;
  \item \textbf{Fase III (Monitoraggio strutturale):} due linee di analisi parallele si occupano
        rispettivamente del monitoraggio metrico aggregato, tramite \textit{threshold}, z-score e
        Isolation Forest, e del monitoraggio comportamentale sul \textit{Path Adherence Score}
        tramite \gls{CUSUM}. Entrambe convergono in un validatore strutturale finale;
  \item \textbf{Fase IV (Sintesi e \textit{Alert}):} l'aggregazione dei risultati produce un
        \textit{alert} strutturato che indica proprietà a rischio, \textit{lead time} stimato, la
        causa radice e il \textit{compliance set} impattato.
\end{itemize}

Attraverso l'adozione di questo approccio stratificato, l'output finale del \textit{framework}
rappresenta a tutti gli effetti un giudizio di \textit{compliance} predittivo completo e non un
semplice segnale di anomalia statistica. Nelle sezioni seguenti vengono formalizzate le strutture
di rappresentazione (Sezione~\ref{sec:rappresentazione_sistema}) e i meccanismi algoritmici della
\textit{pipeline} (Sezione~\ref{sec:pipeline_forecasting}).

\begin{figure}[H]
	\centering
	\includegraphics[width = 0.9\textwidth, height=0.45\textheight,trim={0 0 0 0}, clip]{images/framework_overview}
	\caption[Struttura del framework]{Struttura del framework}
	\label{fig:framework_overview}
\end{figure}

% =========================================================
\section{Rappresentazione del Sistema}
\label{sec:rappresentazione_sistema}
% =========================================================

I Grafi Temporali Attribuiti (\gls{ATG}) permettono di descrivere con precisione il comportamento
di un sistema distribuito a \textit{runtime}. Essi catturano infatti sia le metriche interne dei
microservizi sui nodi, sia le metriche relative alle interazioni sugli archi, consentendo di
costruire modelli di \textit{forecasting} su ciascuna di queste serie temporali. Nonostante ciò,
nel momento in cui il sistema è soggetto a molteplici proprietà certificate, l'utilizzo esclusivo
di un grafo temporale attribuito presenta un limite strutturale significativo; il grafo non permette
infatti di esprimere nativamente la semantica logica che lega un insieme di componenti a una
specifica proprietà certificata.
Questa criticità si accentua ulteriormente in presenza di componenti condivisi tra più percorsi
certificati. Si consideri, ad esempio, una generica risorsa condivisa tra un flusso critico ad alta
priorità e un flusso secondario. Un modello basato esclusivamente su \gls{ATG} può rilevare un
aumento della latenza lungo il flusso critico e correlarlo alla saturazione della risorsa, ma non
dispone dell'informazione semantica per dedurre che il carico del flusso secondario costituisce la
reale causa originaria. In uno scenario di questo genere, l'identificazione della causa radice
richiede un'analisi empirica sui dati storici. Tuttavia, tale approccio presenta due limiti
significativi. Da un lato può infatti introdurre costi computazionali elevati, fino a quadratici in
approcci basati su correlazione a coppie su tutte le metriche disponibili, dall'altro comporta il
rischio che eventuali interferenze tra diverse proprietà rimangano invisibili, finché non si
manifestano esplicitamente nei dati. Per superare questi limiti, viene adottato l'approccio ibrido
a tre \textit{layer} introdotto ad inizio capitolo. I tre \textit{layer} operano su scale temporali
differenti e svolgono funzioni complementari ma non sovrapponibili. L'ipergrafo di certificazione è
definito a \textit{design-time}; viene infatti determinato in fase di configurazione del sistema e
subisce modifiche solo in caso di cambiamenti architetturali o variazioni nelle proprietà
certificate. Al contrario, l'\gls{ATG} è una struttura a \textit{runtime} che evolve costantemente
in base al ritmo della telemetria. Infine, la funzione di \textit{mapping} rappresenta il ponte tra
i primi due livelli e consente al sistema di tradurre le previsioni metriche in valutazioni di
\textit{compliance}, mantenendo intatto il valore semantico della certificazione durante tutto il
ciclo di vita.

% ---------------------------------------------------------
\subsection{Ipergrafo di Certificazione e Compliance Sets}
\label{sec:ipergrafo_certificazione}
% ---------------------------------------------------------

Sulla base dei limiti analizzati, emerge la necessità di definire un formalismo che sia in grado di
rappresentare in maniera esplicita i vincoli di conformità tra molteplici componenti. Viene quindi
introdotta la nozione di \textit{compliance set} intesa come l'insieme minimale di componenti la cui
conformità contemporanea configura una condizione necessaria affinché una proprietà certificata sia
soddisfatta; quindi, un \textit{compliance set} non rappresenta un semplice percorso nel grafo,
bensì costituisce un vincolo logico di tipo AND tra tutti i componenti in esso contenuti. La
violazione dello \gls{SLA} da parte di un solo elemento fa sì che l'intera proprietà venga violata.
Questa semantica non emerge dalla struttura arco-nodo di un grafo tradizionale, in quanto
quest'ultimo modella solo le relazioni binarie tra coppie di nodi. Un \textit{compliance set}
invece può coinvolgere un numero arbitrario di elementi in una singola relazione atomica. \\
Il formalismo degli ipergrafi permette, attraverso una struttura matematica appropriata, di
rappresentare questa semantica. Un ipergrafo rappresenta infatti una generalizzazione del grafo
standard in cui un iperarco può connettere un insieme arbitrario di nodi; un iperarco non è quindi
vincolato ad avere esattamente due estremi, bensì può contenere un qualsiasi numero finito di nodi.
Formalmente, un ipergrafo è una coppia $H = (V, E_H)$\footnote{La notazione
originale di Berge adotta $X$ come simbolo per l'insieme dei vertici, definendo l'ipergrafo come
la coppia $H = (X, E)$. L'adattazione a $V$, utilizzata in questo lavoro, è coerente con la
notazione standard dei grafi adottata nel resto del \textit{framework} e non altera la definizione
sostanziale.} ~\cite{berge1989}, dove:

\begin{itemize}
  \item $V$ è l'insieme dei nodi rappresentanti i microservizi;
  \item $E_H$ è la famiglia degli iperarchi, dove ogni iperarco $h \in E_H$ è un sottoinsieme non
        vuoto di $V$.
\end{itemize}

Nel contesto del \textit{framework}, ogni proprietà certificata viene quindi mappata su un
iperarco dell'ipergrafo di certificazione. L'utilizzo degli ipergrafi, inoltre, permette di
visualizzare immediatamente le interferenze strutturali; i nodi condivisi tra più iperarchi 
indicano infatti componenti critici in cui un singolo degrado prestazionale può innescare violazioni a cascata su diverse certificazioni.

\subsubsection{Definizione Formale}
\label{sec:def_formale_ipergrafo}

A ogni proprietà $\Phi_i$ che il sistema distribuito deve garantire, viene associato un iperarco
$H_{\Phi_i} \subseteq V$ che costituisce il \textit{compliance set} della proprietà. La
semantica dell'iperarco risponde a una logica AND atomica: la proprietà $\Phi_i$ è considerata
verificata se, e solo se, ogni singolo componente appartenente all'insieme rispetta il proprio
\gls{SLA} individuale. Formalmente, tale relazione è espressa come:
\begin{equation}
\label{eq:compliance_condition}
\Phi_i \text{ verificata} \iff \forall\, v \in H_{\Phi_i} : \text{SLA}_v \text{ rispettato}
\end{equation}
In questa espressione, $\text{SLA}_v$ rappresenta la specifica degli \gls{SLA} del componente $v$.
La violazione dell'equazione~\eqref{eq:compliance_condition} per almeno un nodo appartenente a $H_{\Phi_i}$ 
implica la violazione di $\Phi_i$ sull'intero \textit{compliance set}. Questa asimmetria tra la rigidità 
del requisito di conformità universale e la facilità di violazione rende l'intero vincolo 
vulnerabile al degrado di una singola parte. L'ipergrafo di certificazione del sistema è la struttura 
che aggrega tutti i \textit{compliance sets} delle proprietà monitorate:
\begin{equation}
\label{eq:h_cert}
H_{\text{cert}} = \bigl(V,\; \{H_{\Phi_1}, H_{\Phi_2}, \ldots, H_{\Phi_m}\}\bigr)
\end{equation}
$H_{\text{cert}}$ rappresenta una struttura statica che descrive le dipendenze logiche definite in
fase di progettazione architetturale e viene aggiornata solo quando cambiano le proprietà
certificate o l'architettura del sistema. $V$ rappresenta l'insieme dei microservizi.

\subsubsection{Istanziazione sul Sistema di Telemedicina}
\label{sec:istanziazione_telemedicina}

Nel sistema di telemedicina critica introdotto nel Capitolo~\ref{chap:stato_arte}, l'insieme dei
microservizi è:
\begin{equation}
\label{eq:V_telemed}
V = \{\text{SGW},\, \text{SIN},\, \text{SDB},\, \text{SAL},\, \text{SBL}\}
\end{equation}
dove le sigle identificano rispettivamente l'\textit{API Gateway} (SGW), il \textit{Data Ingestion
Service} (SIN), il \textit{Database Service} (SDB), l'\textit{Alert Service} (SAL) e il
\textit{Billing Service} (SBL). Il sistema è soggetto a due proprietà certificate con requisiti
distinti.
La proprietà $\Phi_{\text{crit}}$ (\textit{Real-time Dependability}) impone che la latenza
end-to-end del flusso clinico rimanga inferiore a 500\,ms e che la disponibilità del percorso
superi il 99,99\% su base mensile. Il \textit{compliance set} associato a questa proprietà include
tutti i servizi che costituiscono il flusso clinico:
\begin{equation}
\label{eq:h_crit}
H_{\text{crit}} = \{\text{SGW},\, \text{SIN},\, \text{SDB},\, \text{SAL}\}
\end{equation}
La proprietà $\Phi_{\text{adm}}$ (\textit{Billing Compliance}) richiede la corretta esecuzione di
ogni operazione di fatturazione e la coerenza della base documentale. Il relativo
\textit{compliance set} include i servizi coinvolti nel flusso amministrativo:
\begin{equation}
\label{eq:h_adm}
H_{\text{adm}} = \{\text{SGW},\, \text{SBL},\, \text{SDB}\}
\end{equation}
\subsubsection{Funzione Shared e Identificazione delle Interferenze}
\label{sec:funzione_shared}

Il valore architetturale dell'ipergrafo di certificazione emerge con particolare chiarezza quando
si analizza la relazione tra \textit{compliance sets} distinti. Per ogni coppia di proprietà
$(\Phi_i, \Phi_j)$ con $i \neq j$, si definisce la funzione \textit{Shared} che identifica i
componenti condivisi tra i rispettivi \textit{compliance sets}:
\begin{equation}
\label{eq:shared}
\text{Shared}(H_{\Phi_i}, H_{\Phi_j}) = H_{\Phi_i} \cap H_{\Phi_j}
\end{equation}
Ogni componente $v \in \text{Shared}(H_{\Phi_i}, H_{\Phi_j})$ rappresenta un \textbf{punto di
interferenza strutturale}. In tali nodi, il degrado prestazionale può propagarsi simultaneamente
verso entrambe le proprietà; allo stesso modo, il carico generato dalle attività afferenti a una
proprietà può deteriorare la qualità del servizio dell'altra.
Applicando la funzione al caso di studio della telemedicina, si ottiene:
\begin{equation}
\label{eq:shared_telemed}
\text{Shared}(H_{\text{crit}}, H_{\text{adm}}) = H_{\text{crit}} \cap H_{\text{adm}}
= \{\text{SGW},\, \text{SDB}\}
\end{equation}
Il \textit{gateway} SGW e il \textit{database} SDB sono elementi condivisi tra le due proprietà.
In particolare, il servizio SDB emerge come il nodo di interferenza critico. Un incremento
sostenuto del traffico legato al flusso amministrativo può infatti saturare il \textit{connection
pool} di SDB, degradando la latenza con cui SIN accede alla persistenza. In questo modo, la
proprietà $\Phi_{\text{crit}}$ viene minacciata attraverso una catena causale che si origina
interamente al di fuori del \textit{compliance set} del flusso clinico. SGW, pur appartenendo a
$\text{Shared}(H_{\text{crit}}, H_{\text{adm}})$, non genera un componente di interferenza
esplicito in $M^{\text{interf}}$: essendo il \textit{gateway} d'ingresso del sistema, non riceve
archi interni da altri nodi di $V$ che convoglino traffico quantificabile verso di esso.
L'interferenza su SGW, se presente, si manifesterebbe come saturazione di $\text{cpu}_{\text{SGW}}$
o $\text{mem}_{\text{SGW}}$, già incluse in $M^{\text{direct}}$
(Sezione~\ref{sec:feature_selection}).
È fondamentale sottolineare che tale informazione possiede natura topologica e non statistica; essa
emerge infatti dall'intersezione di due insiemi definiti a livello logico e non richiede
l'elaborazione di serie storiche o correlazioni tra metriche di percorsi differenti. Mentre un
sistema di monitoraggio basato esclusivamente su \gls{ATG} rileverebbe tale interferenza solo a
seguito di una manifestazione empirica nei dati, l'ipergrafo la rende esplicitamente rappresentabile
già in fase di progettazione.

\subsubsection{Proprietà Strutturali}
\label{sec:proprieta_strutturali}

L'ipergrafo $H_{\text{cert}}$~\eqref{eq:h_cert} soddisfa alcune proprietà strutturali determinanti per la
correttezza e l'efficienza computazionale del \textit{framework}. La \textit{completezza} richiede
che ogni componente $v \in V$ che partecipa attivamente al sistema appartenga ad almeno un
iperarco; difatti, qualora un nodo non fosse incluso in alcun \textit{compliance set}, il suo
monitoraggio non troverebbe giustificazione all'interno del \textit{framework} in quanto il suo
stato non concorrerebbe al soddisfacimento di alcuna proprietà. La \textit{relazione di dominanza
tra proprietà}, invece, analizza le relazioni di inclusione tra gli iperarchi. Se
$H_{\Phi_i} \subseteq H_{\Phi_j}$, allora la proprietà $\Phi_j$ risulta logicamente più
restrittiva, in quanto coinvolge un numero maggiore di componenti il cui \gls{SLA} deve essere
simultaneamente soddisfatto. Ogni violazione di $\Phi_i$ implica matematicamente anche la
violazione di $\Phi_j$, poiché il componente guasto appartenente a $H_{\Phi_i}$ appartiene per
definizione anche all'insieme dominante $H_{\Phi_j}$. Questo ordine parziale può essere utilizzato 
per ottimizzare la sequenza di valutazione. Verificare  prima i \textit{compliance sets} più restrittivi 
consente, qualora risultino conformi, di inferire automaticamente la conformità dei set meno restrittivi 
contenuti in essi. In questo modo si riduce il numero di controlli necessari nei cicli operativi in cui il sistema è sano.
È necessario osservare che la proprietà di completezza è definita rispetto a $V$, garantendo che
ogni componente monitorato contribuisca ad almeno una proprietà certificata, senza però vincolare
la composizione di $V$ stessa. Di conseguenza, il perimetro del \textit{framework} coincide
esattamente con i componenti inclusi esplicitamente in $V$. Componenti infrastrutturali trasversali,
come \textit{load balancer} o \textit{DNS resolver}, il cui fallimento può compromettere più
proprietà simultaneamente, devono essere inclusi esplicitamente in $V$ e nei rispettivi iperarchi
affinché il loro impatto venga modellato correttamente. Al contrario, componenti esclusi da $V$
costituiscono invece un \textit{blind spot} del \textit{framework} e il loro effetto sulla
\textit{compliance} non è rappresentato dal modello. Questa criticità va quindi considerata come
una limitazione della fase di modellazione, piuttosto che un limite intrinseco del sistema reale.

Un vantaggio dell'ipergrafo risiede nella sua \emph{modularità} rispetto all'evoluzione del
sistema. Difatti, quando viene introdotta una nuova proprietà certificata $\Phi_{m+1}$ è
sufficiente aggiungere il relativo iperarco $H_{\Phi_{m+1}}$ senza la necessità di dover
modificare i \textit{compliance sets} esistenti. Analogamente, quando un componente viene aggiunto
all'architettura esso viene inserito nell'iperarco corrispondente, operazione svolta in fase di
aggiornamento del modello di certificazione, mantenendo inalterata la struttura degli altri.
Questa flessibilità è fondamentale in sistemi caratterizzati da architetture in evoluzione continua;
difatti, l'aggiunta di una nuova proprietà non necessita di ridefinire tutte le strutture di
monitoraggio esistenti.

% ---------------------------------------------------------
\subsection{Grafo Temporale Attribuito e Metriche Osservabili}
\label{sec:atg_metriche}
% ---------------------------------------------------------

Il Layer~2 dell'architettura ibrida è rappresentato dall'\gls{ATG}; quest'ultimo modella
l'evoluzione dinamica del sistema distribuito attraverso serie temporali di metriche osservabili.
A differenza del Layer~1, il quale definisce le proprietà certificate, il Layer~2 descrive il
comportamento del sistema a \textit{runtime}, attraverso le misurazioni raccolte dall'infrastruttura
di telemetria.

\subsubsection{Definizione Formale}
\label{sec:def_formale_atg}
Un Grafo Temporale Attribuito al tempo $t$ è definito dalla quadrupla:
\begin{equation}
\label{eq:atg_def}
G_t = (V,\; E_t,\; \mathbf{X}_{V,t},\; \mathbf{X}_{E,t})
\end{equation}
I quattro elementi della definizione hanno ruoli distinti e complementari.

\begin{itemize}
  \item $\mathbf{V}$: rappresenta l'insieme dei nodi, ossia i microservizi o componenti
        strutturali che costituiscono il sistema distribuito. $V$ è condiviso tra Layer~1 e
        Layer~2; difatti, gli stessi componenti che compaiono nei \textit{compliance sets}
        dell'ipergrafo compaiono anche come nodi nell'\gls{ATG}. Questa condivisione costituisce
        il presupposto necessario per la funzione di \textit{mapping} descritta nella
        Sezione~\ref{sec:funzione_mapping};
  \item $\mathbf{E_t \subseteq V \times V}$: indica l'insieme degli archi attivi al tempo $t$.
        Un arco $(u,v) \in E_t$ rappresenta una chiamata di servizio tra il componente $u$ e il
        componente $v$ nell'intervallo temporale corrispondente a $t$. A differenza dei nodi,
        l'insieme degli archi è variabile nel tempo. Nuovi archi compaiono infatti quando vengono
        attivate nuove dipendenze comunicative, mentre archi esistenti scompaiono quando i
        \textit{circuit breaker} interrompono percorsi degradati o quando il sistema viene
        riconfigurato;
  \item $\mathbf{X_{V,t} \in \mathbb{R}^{|V| \times d_V}}$: costituisce la matrice delle
        \textit{feature} di nodo al tempo $t$. Ogni riga $i$ contiene il vettore relativo allo
        stato interno del componente $v_i$. Queste \textit{feature} sono persistenti, ovvero
        assumono valori definiti anche in assenza di traffico in transito sul nodo. Un
        microservizio in stato di \textit{idle}, ad esempio, ha un utilizzo della \gls{CPU}
        prossimo allo zero e una dimensione stabile dell'\textit{heap} di memoria. Questi dati
        sono osservabili e informativi anche senza che il nodo stia elaborando richieste. La
        struttura duale del grafo è compatibile con i principali approcci di \textit{representation
        learning} su grafi dinamici, tra cui i \gls{TGN}~\cite{rossi2020}, che adottano
        un'architettura analoga, distinguendo tra uno stato persistente per nodo
        (\textit{memory module}) e \textit{feature} effimere associate agli eventi di interazione;
  \item $\mathbf{X_{E,t} \in \mathbb{R}^{|E_t| \times d_E}}$: rappresenta la matrice delle
        \textit{feature} di arco al tempo $t$. Ogni riga $j$ contiene il vettore associato
        all'interazione $e_j \in E_t$. A differenza delle \textit{feature} di nodo, le
        \textit{feature} di arco sono effimere: esse esistono solo in presenza di traffico
        sull'arco. Parametri come la latenza di una chiamata, il tasso di errore e il
        \textit{throughput} sono proprietà che nascono e finiscono con l'interazione, e quindi
        legate alla durata dell'interazione stessa.
\end{itemize}

\subsubsection{Ruolo Diagnostico delle Feature}
\label{sec:ruolo_diagnostico}

Le \textbf{\textit{feature} di nodo} rappresentano la sorgente primaria per l'individuazione dei
segnali di degrado lento, definiti come \textit{compliance drift}. Un \textit{memory leak}, ad
esempio, si manifesta attraverso una crescita graduale di $\text{mem}_v(t)$ nell'arco di ore o
giorni, molto prima che la latenza degli archi incidenti superi le soglie di allarme. Analogamente,
la saturazione progressiva del \textit{connection pool} di un \textit{database} è visibile su
$\text{pool}_v(t)$ minuti prima che le richieste inizino a mettersi in coda e la latenza esploda.
In questi scenari, il mancato monitoraggio delle \textit{feature} di nodo fa sì che vengano
prodotti \textit{alert} con un \textit{lead time} insufficiente per un intervento preventivo
efficace. Le \textbf{\textit{feature} di arco}, al contrario, rappresentano la superficie
osservabile delle interazioni e costituiscono la metrica di \textit{compliance} più diretta per
requisiti di latenza o disponibilità. La violazione di uno \gls{SLA} di latenza si manifesta
direttamente come $L_e(t) > L_{\max}$ su un arco del percorso certificato. Il \textit{forecasting}
delle \textit{feature} di arco permette di stimare la latenza futura che viene poi aggregata per
ottenere la previsione della proprietà \textit{end-to-end}.

Nel sistema di telemedicina, le \textit{feature} di nodo associate a ciascun microservizio
$v \in V$ sono raggruppate nel vettore:
\begin{equation}
\label{eq:node_features}
\mathbf{x}_v(t) = \bigl[\text{cpu}_v(t),\; \text{mem}_v(t),\; \text{pool}_v(t),\; \text{gc}_v(t)\bigr]
\end{equation}
Il vettore è definito in modo uniforme su tutti i componenti del sistema, rappresentando lo spazio
delle \textit{feature} disponibili nel Layer~2. La selezione operativa delle metriche rilevanti per
il \textit{forecasting}, che può escludere alcune dimensioni del vettore per specifici nodi, è
demandata alla \textit{feature selection} della Fase~I (Sezione~\ref{sec:feature_selection}). \\
In questa definizione, $\text{cpu}_v(t)$ indica l'utilizzo della \gls{CPU} in percentuale. 
Il parametro $\text{mem}_v(t)$ denota invece l'occupazione della memoria principale del processo in MB che, 
nel sistema di riferimento basato su \gls{JVM}, corrisponde all'occupazione dell'\textit{heap}; 
$\text{gc}_v(t)$ esprime la durata delle pause di gestione della memoria in millisecondi e, 
nel sistema di riferimento, corrisponde alle pause del \textit{garbage collector} \gls{JVM}.
La metrica $\text{pool}_v(t)$ rappresenta invece il tasso di saturazione delle connessioni. 
Il suo significato specifico dipende dal ruolo del nodo: per i microservizi applicativi come 
SIN misura il \textit{connection pool outbound} verso il \textit{database}, 
mentre per il database SDB indica il \textit{pool incoming} delle connessioni dai \textit{client}. 
Per SIN, il \textit{pool} in uscita verso SDB non viene incluso direttamente in $M^{\text{direct}}$ 
(Sezione~\ref{sec:feature_selection}) poiché, nella configurazione del sistema di riferimento, 
la sua saturazione costituisce tipicamente una conseguenza diretta della saturazione del \textit{pool incoming} 
di SDB lato \textit{server}, il quale rappresenta il vero segnale \textit{root-cause}. 
La rilevanza di queste metriche varia a
seconda del componente. Per SGW e SAL, che non gestiscono \textit{pool} di connessioni, il
contributo più significativo proviene da \gls{CPU} e memoria, mentre $\text{pool}_v(t)$ è la
metrica determinante per il \textit{database} SDB; infine $\text{gc}_v(t)$ segnala il degrado
nelle fasi di allocazione per SIN, che esegue normalizzazione dati su \gls{JVM}. 
Per quanto riguarda le interazioni $e = (u,v) \in E_t$, le \textit{feature} di arco sono definite dal vettore:
\begin{equation}
\label{eq:edge_features}
\mathbf{x}_e(t) = \bigl[L_e(t),\; \varepsilon_e(t),\; \theta_e(t)\bigr]
\end{equation}
dove $L_e(t)$ indica la latenza media delle chiamate sull'arco $e$ in ms, $\varepsilon_e(t)$ è il
tasso di errore come frazione delle richieste fallite, e $\theta_e(t)$ è il \textit{throughput} in
richieste al secondo.
Questa modellazione duale consente al \textit{framework} di intercettare il degrado interno prima
che questo si propaghi all'esterno, trasformando segnali deboli a livello di nodo in previsioni
azionabili sulla \textit{compliance} globale.

I due scenari seguenti illustrano concretamente la complementarità tra \textit{feature} di nodo e
\textit{feature} di arco. Per semplicità illustrativa, si assume che il \textit{budget} totale di
500\,ms imposto dalla proprietà \emph{Real-time Dependability} venga ripartito allocando soglie
locali indicative di 250\,ms per i singoli archi critici. Entrambe le tabelle riportano valori a
titolo puramente illustrativo.

\begin{table}[htbp]
\centering
\caption{Scenario \textit{memory leak} in SIN: le metriche di nodo segnalano il degrado con
largo anticipo rispetto alla metrica di arco (valori illustrativi).}
\label{tab:scenario_memory_leak}
\begin{tabular}{ccccc}
\toprule
\textbf{Giorno} & $\boldsymbol{L_{e_1}}$\textbf{(ms)} & $\boldsymbol{L_{e_2}}$\textbf{(ms)} & \textbf{mem\_SIN (MB)} & \textbf{cpu\_SIN (\%)} \\
\midrule
0  & 45 & 80  & 512  & 28 \\
5  & 45 & 158 & 972  & 44 \\
10 & 46 & 268 & 1507 & 63 \\
\midrule
\multicolumn{5}{c}{\footnotesize\textbf{Soglia \gls{SLA}:} $L_{e_1} \leq 250\;\text{ms}$, $L_{e_2} \leq 250\;\text{ms}$ (\textit{budget} ripartito)} \\
\bottomrule
\end{tabular}
\end{table}

La Tabella~\ref{tab:scenario_memory_leak} mostra come la memoria del servizio SIN subisca un
incremento del 194\% in dieci giorni, mentre la latenza dell'arco $e_2$ supera la soglia definita
dallo \gls{SLA} solo al decimo giorno. Un sistema che monitora esclusivamente le \textit{feature}
di arco genererebbe l'alert al decimo giorno senza alcun preavviso utile. Al contrario, un sistema
che monitora anche le \textit{feature} di nodo può prevedere la violazione già al quinto giorno,
quando $\text{mem}_{\text{SIN}}$ supera i 900\,MB con un \textit{trend} di crescita con tasso
accelerato. In questo scenario, il coefficiente di correlazione tra $\text{mem\_SIN}$ e $L_{e_2}$
è prossimo a 0,99, risultato empiricamente coerente con una relazione causale, la cui direzione
viene verificata nella Fase~II tramite il test di Granger.

\begin{table}[htbp]
\centering
\caption{Scenario interferenza su SDB: la saturazione del \textit{connection pool} precede la
violazione di latenza sul percorso critico (valori illustrativi).}
\label{tab:scenario_interferenza}
\begin{tabular}{ccccc}
\toprule
\textbf{Tempo} & $\boldsymbol{L_{e_2}}$\textbf{(ms)} & $\boldsymbol{\theta_{e_5}}$\textbf{(req/s)} & \textbf{pool\_SDB (\%)} & \textbf{Stato} \\
\midrule
08:00 & 82  & 5  & 38 & Normale \\
08:30 & 127 & 25 & 81 & In degrado \\
09:00 & 450 & 29 & 98 & \textit{Failure} \\
\midrule
\multicolumn{5}{c}{\footnotesize\textbf{Soglia \gls{SLA}:} $L_{e_2} \leq 250\;\text{ms}$;\; soglia \textit{pool}: 90\%} \\
\bottomrule
\end{tabular}
\end{table}

La Tabella~\ref{tab:scenario_interferenza} illustra un caso di interferenza \textit{cross-property}.
Il \textit{throughput} del flusso amministrativo $\theta_{e_5}$ aumenta da 5 a 29\,req/s in
un'ora, saturando progressivamente il \textit{connection pool} del \textit{database} SDB. Alle
08:30 la saturazione del \textit{pool} raggiunge l'81\%, avvicinandosi pericolosamente alla soglia
critica del 90\%, mentre la latenza $L_{e_2}$ rimane conforme. Un sistema di monitoraggio basato
solo su \textit{feature} di arco, senza quindi accesso alla \textit{feature} di nodo
$\text{pool}_{\text{SDB}}$, non disporrebbe dei dati necessari per generare un preavviso utile in
questo momento. La \textit{feature} di nodo, in questo caso, rende visibile il degrado circa trenta
minuti prima della violazione esplicita. Questo anticipo costituisce il \textit{lead time}
disponibile per l'attuazione di un intervento preventivo.

\paragraph{Selezione delle Metriche}
\label{sec:selezione_metriche}

Il formalismo dell'\gls{ATG} non impone un \textit{set} predefinito di metriche. La distinzione
tra \textbf{\textit{feature} di nodo} e \textbf{\textit{feature} di arco} ha una valenza
concettuale e non prescrittiva. Di conseguenza, in un'applicazione reale, le annotazioni del grafo
temporale dipendono dal dominio specifico e dallo \textit{stack} tecnologico in uso. Per quanto
riguarda l'analisi dello stato dei nodi, possono essere utilizzate metriche relative a: la
saturazione delle risorse, come \gls{CPU} e memoria, contatori applicativi, come numero di
richieste in coda e \textit{cache hit rate} oppure eventi di gestione runtime, come pause del
\textit{garbage collector} e saturazione dei \textit{thread pool}. Per quanto riguarda gli archi,
invece, tipicamente vengono utilizzati indicatori di qualità delle interazioni; tra questi si hanno
la latenza \textit{end-to-end}, il \textit{throughput}, il tasso di errore.
La scelta delle metriche da utilizzare dipende strettamente dalle proprietà che devono essere
monitorate. Nel caso di proprietà come la \textit{Real-time Dependability} vengono utilizzate
metriche relative ai parametri temporali, come ad esempio latenza e \textit{jitter}. Diversamente,
per proprietà come la \textit{Billing Compliance}, si utilizzano invece metriche transazionali;
tra queste figurano il \textit{throughput} per quantificare i volumi di fatturazione ($\theta_e$)
e il tasso di errore ($\varepsilon_e$) per intercettare transazioni fallite o scartate. Questa
flessibilità fa sì che il \textit{framework} proposto sia applicabile a qualsiasi sistema
distribuito, purché questo sia dotato di telemetria strutturata secondo la distinzione tra stato
dei nodi e qualità delle interazioni indipendentemente dalle tecnologie utilizzate.

\paragraph{Integrazione con l'Infrastruttura di Osservabilità}
\label{sec:integrazione_infrastruttura}

Il \textit{framework} è progettato per integrarsi con le \textit{pipeline} di osservabilità
standard dell'ecosistema \textit{cloud-native}. Per l'instrumentazione, strumenti come
OpenTelemetry\footnote{\url{https://opentelemetry.io/}} o le librerie \textit{client} di
Prometheus\footnote{\url{https://prometheus.io/}} permettono di produrre nativamente le serie
temporali che alimentano il Layer~2.
Per quanto riguarda la persistenza delle serie temporali, essa può essere affidata a soluzioni
come Prometheus TSDB o TimescaleDB\footnote{\url{https://www.timescale.com/}}; queste soluzioni
permettono di garantire la capacità di gestione di dati ad alta frequenza, necessaria per le
finestre di analisi del \textit{framework}. Per la ricostruzione automatica della topologia $E_t$,
invece, strumenti di \textit{distributed tracing} come Jaeger\footnote{\url{https://www.jaegertracing.io/}}
e Zipkin\footnote{\url{https://zipkin.io/}} permettono di ricostruire le dipendenze comunicative
attive a partire dagli \textit{span} delle richieste.\\
L'architettura si mantiene quindi agnostica rispetto allo \textit{stack} tecnologico specifico;
l'unico requisito è costituito dalla capacità di esportare le serie temporali in un formato
strutturato, come CSV o Parquet, per l'ingestione nella \textit{pipeline} di \textit{forecasting}.

\subsubsection{Il Probabilistic Behavioral Overlay}
\label{sec:pbo}

Al fine di quantificare con precisione fenomeni quali latenza o saturazione delle risorse vengono
utilizzate le metriche associate ai nodi e agli archi dell'\gls{ATG}; tuttavia, questi dati non
sono sufficienti per stabilire se il sistema sia conforme al comportamento atteso oppure se stia
progressivamente deviando verso traiettorie impreviste. Per superare questo limite, il
\textit{framework} sfrutta un livello di astrazione superiore, denominato \textbf{\gls{PBO}}.
Questo livello permette al \textit{framework} stesso di interpretare le osservazioni quantitative
sotto forma di comportamento globale del sistema. È importante però distinguere \gls{ATG} e
\gls{PBO}, in quanto offrono due viste complementari ma distinte del sistema.

L'\gls{ATG} cattura una rappresentazione \emph{fine-grained}; difatti, per ogni istante temporale,
vengono registrati i valori puntuali delle metriche su nodi e archi attivi. Questa granularità
risulta fondamentale sia per la costruzione di modelli di previsione accurati sulle singole serie
temporali, sia per l'identificazione delle cause radice di eventuali anomalie a livello di singolo
componente. Il \gls{PBO} fornisce invece una rappresentazione \textit{coarse-grained} del
comportamento; difatti, anziché descrivere i valori esatti delle metriche, permette di catturare
la frequenza con cui le richieste seguono ciascun percorso disponibile. Tramite il suo utilizzo,
anche in assenza di anomalie nelle metriche, è possibile identificare una deviazione
comportamentale. Ad esempio, una diminuzione progressiva del traffico sul percorso clinico non
necessariamente produce variazioni immediate nelle metriche puntuali degli archi del percorso
primario; tuttavia, questa deviazione può essere identificata attraverso il \gls{PBO}. Da un punto
di vista formale, la costruzione del \gls{PBO} trae ispirazione concettuale dai principi del
\textit{Process Mining} e, in particolare, dalla struttura dati dei \textit{Directly-Follows
Graphs}~\cite{vanderaalst2016}. Tuttavia, mentre nel \textit{process mining}
tradizionale tali strumenti vengono impiegati prevalentemente per l'analisi a posteriori e la
\textit{conformance checking} dei processi, nel presente \textit{framework}
la struttura probabilistica viene riadattata e applicata dinamicamente al monitoraggio continuo e
al \textit{forecasting} delle interazioni nei sistemi a microservizi. Formalmente, il \gls{PBO}
è definito come il grafo:
\begin{equation}
\label{eq:pbo_def}
G_{\text{Behavior}}(t) = (V,\; E_{\text{all}},\; W_t)
\end{equation}
dove:

\begin{itemize}
  \item $\mathbf{V}$: rappresenta il medesimo insieme di nodi dell'\gls{ATG}.
  \item $\mathbf{E_{\text{all}}}$: rappresenta l'insieme di tutti gli archi osservati nel sistema
        sull'intera finestra storica di addestramento. Questo insieme include sia gli archi attivi
        al tempo $t$, sia i percorsi di \textit{fallback} che non risultano necessariamente attivi
        in ogni istante;
  \item $\mathbf{W_t}$: è la matrice dei pesi di transizione stocastici. Il peso di ogni arco
        $(u,v)$ indica la quantità di traffico totale in uscita dal nodo $u$ che sceglie quel
        percorso specifico in un dato intervallo di tempo. La matrice $W_t$ è, per costruzione,
        stocastica per righe limitatamente ai nodi con archi uscenti attivi: la somma dei pesi
        uscenti da ogni nodo non terminale è pari a 1, mentre la riga corrispondente a un nodo
        terminale risulta nulla, in quanto tale nodo non instrada traffico verso altri componenti.
\end{itemize}

\paragraph{Il \textit{Path Adherence Score}}
\label{sec:path_adherence_score}

Il \gls{PBO} permette di definire una metrica che misura quanto il sistema stia effettivamente
seguendo il percorso comportamentale certificato per una data proprietà. In particolare, dato il
percorso certificato $P_{\text{cert}} = (v_1, v_2, \ldots, v_n)$ associato alla proprietà $\Phi_i$,
il \textit{Path Adherence Score} al tempo $t$ è definito come il prodotto dei pesi stocastici lungo
il percorso:
\begin{equation}
\label{eq:pa_score}
PA(P_{\text{cert}},\; t) = \prod_{k=1}^{n-1} w_{v_k, v_{k+1}}(t)
\end{equation}
Il valore di $PA(P_{\text{cert}}, t)$ è compreso tra 0 e 1 e si interpreta, sotto ipotesi di
indipendenza delle transizioni (proprietà di Markov), come la probabilità empirica che una
richiesta attraversi l'intero percorso certificato $P_{\text{cert}}$ nella finestra temporale
corrente. Un valore $PA \approx 1$ rappresenta una piena conformità comportamentale;
significa infatti che il traffico segue quasi esclusivamente il percorso certificato. Invece,
$PA \to 0$ significa che una porzione crescente del traffico sta deviando su
percorsi alternativi; ciò segnala un allontanamento strutturale dal comportamento atteso. Questa
formulazione è rigorosamente valida per percorsi lineari ordinati, ed è il caso del sistema di
telemedicina analizzato, dove il flusso clinico segue la sequenza univoca
SGW$\to$SIN$\to$SDB$\to$SAL. Qualora il \textit{compliance set} $H_{\Phi_i}$ presenti invece una
topologia con rami paralleli o bilanciamento di carico, il \gls{PAS} non è applicabile in questa
forma. In tali contesti, infatti, l'aderenza strutturale viene quantificata globalmente attraverso
la distanza tra le matrici di transizione.

Il \textit{Path Adherence Score} costituisce una serie temporale direttamente utilizzabile per il
\textit{forecasting}. Applicando un modello di previsione a $PA(t)$ è infatti possibile stimare
quanto tempo residuo il sistema abbia prima di abbandonare il percorso certificato, anche in
assenza di violazioni metriche puntuali già manifeste.
Per il sistema di telemedicina analizzato, il percorso certificato del flusso clinico corrisponde
alla sequenza $P_{\text{crit}} = (\text{SGW}, \text{SIN}, \text{SDB}, \text{SAL})$, che attraversa
i tre archi della catena clinica in ordine. Il \textit{Path Adherence Score} si ottiene applicando
la formula~\eqref{eq:pa_score} ai pesi degli archi di questo percorso nel \gls{PBO}. Un valore
prossimo a 1 indica che la quasi totalità del traffico clinico scorre lungo il percorso atteso.
Al contrario, una diminuzione del valore segnala un allontanamento strutturale che solitamente
indica e precede l'attivazione di meccanismi di \textit{failover} o la comparsa di errori visibili
sulle metriche di arco.

\paragraph{Calibrazione del Gold Standard e Graph Edit Distance}
\label{sec:gold_standard_ged}

Il \gls{PBO} osservato $G_{\text{Behavior}}(t)$ è confrontato con un \gls{PBO} di riferimento
denominato \textit{PBO gold standard}:
\begin{equation}
\label{eq:pbo_gold}
G_{\text{Behavior}}^{\text{gold}} = (V,\; E_{\text{all}},\; W^{\text{gold}})
\end{equation}
In questa definizione, $E_{\text{all}}$ rappresenta l'insieme completo degli archi osservati nella
finestra storica di addestramento. $W^{\text{gold}}$ identifica invece la matrice dei pesi di
transizione osservata durante un periodo di normale operatività su tutti gli archi di $E_{\text{all}}$.
In questa matrice, i pesi degli archi certificati $E_{\text{cert}}$ risulteranno tipicamente
elevati, mentre quelli dei percorsi di \textit{fallback} saranno inferiori ma non necessariamente
nulli, riflettendo quindi il traffico residuo osservato durante la calibrazione. Il sottoinsieme
$E_{\text{cert}}$ degli archi certificati rimane rilevante per il calcolo del \textit{Path
Adherence Score}, ma non rappresenta l'insieme di archi del \textit{gold standard}. La definizione
del \textit{gold standard} richiede una fase iniziale di calibrazione, durante la quale il sistema
deve operare in condizioni di normale operatività e le distribuzioni del traffico devono
corrispondere fedelmente ai modelli attesi. La conformità di queste condizioni viene verificata
tramite un'analisi retrospettiva sulle metriche di arco e di nodo: le finestre temporali in cui
emergono anomalie nei dati storici vengono escluse dal calcolo di $W^{\text{gold}}$, e il periodo
di osservazione viene esteso fino a ottenere un campione pulito di dimensione sufficiente. La
qualità di questa fase è fondamentale. Un \textit{dataset} di calibrazione contenente anomalie
produrrebbe un \textit{gold standard} inquinato e ciò porterebbe il sistema a generare falsi
negativi durante la fase operativa, andando a vanificare l'intera architettura di rilevamento.
Per questa ragione, il \textit{gold standard} deve essere aggiornato ogni qualvolta il sistema
subisca modifiche strutturali: migrazioni infrastrutturali, variazioni significative nelle
dinamiche di traffico, o qualsiasi altro cambiamento che alteri le distribuzioni attese. Nelle fasi
di transizione architetturale si crea un lasso temporale intermedio in cui il vecchio
\textit{gold standard} è invalidato ma il nuovo non è ancora calibrato. In questo caso, il sistema
gestisce questa finestra di vulnerabilità adottando soglie di tolleranza più ampie sulla distanza
strutturale, contenendo così la generazione di \textit{alert} impropri fino al completamento della
ricalibratura.

In scenari dove l'insieme dei nodi $V$ e degli archi $E_{\text{all}}$ è fisso, la topologia del
grafo non varia nel tempo. In questi casi, lo scostamento strutturale tra il \gls{PBO} osservato e
il \textit{gold standard} si riduce interamente alla variazione dei pesi e può essere quantificato
attraverso la norma di Frobenius $\|W(t) - W^{\text{gold}}\|_F$, senza necessità di ricorrere agli
algoritmi \gls{GED}; difatti, ciò che cambia a ogni istante $t$ sono esclusivamente i pesi $W_t$.
La distanza strutturale viene calcolata attraverso la norma di
Frobenius\footnote{La norma di Frobenius estende la norma euclidea allo spazio 
delle matrici, aggregando in un unico scalare la magnitudo di tutte le 
variazioni elemento per elemento.} tra le matrici dei pesi, la quale
misura la deviazione complessiva come radice della somma dei quadrati degli scarti su tutti gli
archi:
\begin{equation}
\label{eq:frobenius}
\|W(t) - W^{\text{gold}}\|_F = \sqrt{\sum_{u,v} \bigl(w_{uv}(t) - w^{\text{gold}}_{uv}\bigr)^2}
\end{equation}
Il ricorso ad algoritmi di \textit{Graph Edit Distance} (\gls{GED}), definita come il numero
minimo ponderato di operazioni elementari, quali inserimento o cancellazione di nodi e archi e modifica dei
pesi, per trasformare $G_{\text{Behavior}}(t)$ nel \textit{gold standard}, è giustificato nei
contesti in cui la topologia $E_t$ sia effettivamente variabile a \textit{runtime}; ad esempio,
in sistemi con \textit{scaling} dinamico che introduce nuovi archi non presenti nella finestra
storica. Tuttavia, il calcolo esatto della \gls{GED} ha complessità esponenziale nel numero di
nodi dei grafi coinvolti; per questo motivo il \textit{framework} ricorre ad approssimazioni
efficienti basate su metodi derivati dall'algoritmo Ungherese~\cite{riesen2009}, che garantiscono
soluzioni di buona qualità in tempo polinomiale. In scenari dove la topologia è variabile ma il
percorso certificato rimane attivo come sottoinsieme del grafo, un valore elevato di \gls{GED}
combinato con una diminuzione del \textit{Path Adherence Score} costituisce un segnale di allerta
strutturale fondamentale per la classificazione dell'\textit{alert} nella Fase~IV. Questa
combinazione indica che il sistema non solo si sta allontanando dal percorso certificato, ma lo
sta facendo in modo tale da richiedere un numero significativo di trasformazioni per tornare al
comportamento di riferimento. Ciò segnala un potenziale degrado profondo della conformità.

% ---------------------------------------------------------
\subsection{Funzione di Mapping}
\label{sec:funzione_mapping}
% ---------------------------------------------------------

Il \textit{Layer}~3 dell'architettura ibrida è costituito dalla funzione di \textit{mapping}
$\mathcal{M}$, che collega le strutture statiche del \textit{Layer}~1 alle strutture dinamiche
del \textit{Layer}~2. Senza di essa, i due livelli opererebbero in modo indipendente e il
\textit{framework} non potrebbe tradurre le previsioni metriche in giudizi di conformità sulle
proprietà certificate. A partire dalla specifica di ciascun \textit{compliance set}, $\mathcal{M}$
determina in modo deterministico quali metriche del grafo temporale risultino rilevanti per una
data proprietà e quali previsioni locali debbano essere aggregate per ottenere una previsione
globale di \textit{compliance}.

\subsubsection{Definizione Formale}
\label{sec:def_formale_mapping}

La funzione di \textit{mapping} è una relazione bidirezionale tra iperarchi del \textit{Layer}~1
e i sottoinsiemi di archi del \textit{Layer}~2. Nella direzione Layer~1 $\to$ Layer~2, la funzione
associa ogni \textit{compliance set} all'insieme degli archi del grafo temporale il cui degrado
impatta direttamente la proprietà:
\begin{equation}
\label{eq:mapping_def}
\mathcal{M}(H_{\Phi_i}) = A(H_{\Phi_i}) = \{e = (u,v) \in E_t \mid u \in H_{\Phi_i} \wedge v \in H_{\Phi_i}\}
\end{equation}
Nella direzione inversa Layer~2 $\to$ Layer~1, i metadati dell'annotazione semantica permettono di
risalire dalla degradazione di uno specifico arco $e$ all'insieme di proprietà impattate. In
particolare, si ha che: per ogni $H_{\Phi_i}$ tale che $e \in A(H_{\Phi_i})$, la violazione
dell'arco contribuisce alla valutazione di conformità di $\Phi_i$. Questo meccanismo inverso viene
utilizzato nella Fase~IV, dove la funzione di aggregazione traduce le previsioni locali in giudizi
globali di \textit{compliance}. Inoltre, per gli archi esterni identificati come potenziali vettori
di interferenza $M^{\text{interf}}$, il meccanismo inverso sfrutta la funzione \textit{Shared} per
mappare l'anomalia dell'arco di interferenza direttamente alle proprietà che dipendono dal nodo
condiviso saturato. La definizione comprende gli archi interni al \textit{compliance set}, ossia le comunicazioni
dirette tra i componenti del \textit{set}. Invece, gli archi provenienti da nodi esterni, che
convogliano traffico verso componenti condivisi, non rientrano in $A(H_{\Phi_i})$ ma vengono
inclusi nel componente $M^{\text{interf}}$ della \textit{feature selection}
(Sezione~\ref{sec:feature_selection}), ossia il meccanismo dedicato alla gestione delle interferenze
\textit{cross-property}. \\
Nel sistema di telemedicina, la topologia del \textbf{flusso clinico} è definita dalla sequenza
$\text{SGW} \to \text{SIN} \to \text{SDB} \to \text{SAL}$, con gli archi $e_1$, $e_2$ ed $e_3$.
Il \textit{mapping} per la proprietà $\Phi_{\text{crit}}$ risulta quindi:
\begin{equation}
\label{eq:mapping_crit}
A(H_{\text{crit}}) = \{e_1\!:\text{SGW}\!\to\!\text{SIN},\; e_2\!:\text{SIN}\!\to\!\text{SDB},\; e_3\!:\text{SDB}\!\to\!\text{SAL}\}
\end{equation}
Per la proprietà $\Phi_{\text{adm}}$, la topologia del \textbf{flusso amministrativo} segue il
percorso $\text{SGW} \to \text{SBL} \to \text{SDB}$, con gli archi $e_4$ ed $e_5$:
\begin{equation}
\label{eq:mapping_adm}
A(H_{\text{adm}}) = \{e_4\!:\text{SGW}\!\to\!\text{SBL},\; e_5\!:\text{SBL}\!\to\!\text{SDB}\}
\end{equation}
\subsubsection{Strategie di Implementazione del Mapping}
\label{sec:strategie_mapping}

L'adozione di librerie di \textit{machine learning} e algoritmi di analisi su grafi impone un
vincolo di compatibilità: la maggior parte degli strumenti disponibili opera nativamente su grafi
con \textit{feature} sugli archi, non su strutture ipergrafiche. Il Layer~3 deve quindi tradurre
la struttura dell'ipergrafo in una rappresentazione compatibile con il Layer~2, preservando la
semantica dei \textit{compliance sets}. Sono state analizzate due strategie principali.

\textbf{Annotazione Semantica}: mantiene invariata la topologia originale del grafo temporale,
arricchendo ogni arco con metadati che ne dichiarano l'appartenenza agli iperarchi corrispondenti.
Per ogni arco $e = (u,v)$, viene mantenuto l'insieme degli iperarchi a cui appartiene (tutti gli
$H_{\Phi_i}$ tali che $u \in H_{\Phi_i}$ \textbf{e} $v \in H_{\Phi_i}$). Quando il Layer~3
analizza una previsione su un arco $e$, consulta queste annotazioni statiche per determinare quali
proprietà vengono impattate. Questa soluzione non introduce \textit{overhead} a \textit{runtime},
poiché i metadati vengono calcolati una sola volta in fase di definizione del sistema. Questo
meccanismo di annotazione copre gli archi interni ai \textit{compliance sets}. Gli archi che
costituiscono $M^{\text{interf}}$, ossia archi con un solo estremo in $H_{\Phi_i}$, non vengono
annotati con gli iperarchi di cui non sono interni: la loro rilevanza viene determinata
separatamente attraverso la funzione \textit{Shared} durante la \textit{feature selection} della
Fase~I (Sezione~\ref{sec:feature_selection}).

\textbf{Espansione Topologica}: introduce nodi virtuali di sincronizzazione. Per ogni
\textit{compliance set} $H_{\Phi_i}$, si aggiunge al grafo temporale un nodo sintetico $v_{H_i}$
a cui vengono connessi tutti i nodi del \textit{set} tramite archi fittizi. Questa strategia rende
la logica AND esplicita nella topologia, poiché la sua attivazione richiede che tutti i nodi del
\textit{set} siano attivi e conformi. Tuttavia, ciò introduce un \textit{overhead} computazionale
crescente con il numero di proprietà da monitorare; il grafo deve essere infatti arricchito di $m$
nodi, uno per proprietà, le cui \textit{feature} devono essere ricalcolate a ogni passo temporale.

La Tabella~\ref{tab:mapping_strategies} sintetizza i vantaggi e gli svantaggi delle due
metodologie di implementazione.

\begin{table}[htbp]
\centering
\caption{Confronto tra le strategie di implementazione del \textit{mapping}.}
\label{tab:mapping_strategies}
\begin{tabular}{lcc}
\toprule
\textbf{Proprietà} & \textbf{Annotazione Semantica} & \textbf{Espansione Topologica} \\
\midrule
\textbf{Semantica}                   & Implicita (metadati) & AND Esplicita (topologica) \\
\textbf{Topologia invariata}         & Sì                   & No \\
\textbf{\textit{Overhead} spaziale}  & Nullo                & $m$ nodi aggiuntivi \\
\textbf{\textit{Overhead runtime}}   & Nullo                & Ricalcolo \textit{feature} nodi \\
\textbf{Compatibilità algoritmi}     & Tramite \textit{query} su metadati & Nativa \\
\textbf{Identificazione interferenze}& \textit{Query} diretta & Difficile \\
\bottomrule
\end{tabular}
\end{table}

\subsubsection{Scelta dell'Annotazione Semantica}
\label{sec:scelta_annotazione}

Per l'implementazione del \textit{framework} proposto in questo lavoro, viene adottata la
strategia dell'\textbf{Annotazione Semantica} per tre ragioni convergenti. La prima è data
dall'\textit{efficienza computazionale}. In sistemi che monitorano un numero elevato di proprietà,
l'espansione topologica moltiplicherebbe il numero di nodi e renderebbe il grafo più oneroso da
elaborare ad ogni passo temporale; ciò non porterebbe vantaggi al \textit{forecasting} che non
siano ottenibili attraverso le annotazioni. La seconda è la \textit{compatibilità con la telemetria
standard}. Le \textit{pipeline} di osservabilità industriali producono infatti metriche e tracce
associate ad archi e nodi del \textit{service dependency graph} reale, non a nodi virtuali.
L'annotazione semantica consente quindi di arricchire direttamente queste strutture senza richiedere
l'inserimento di nodi virtuali. La terza ragione è la \textit{scalabilità rispetto all'evoluzione
del sistema}. L'aggiunta di una nuova proprietà certificata, attraverso l'utilizzo dell'annotazione
semantica, richiede solo di aggiornare le annotazioni degli archi già preesistenti. Al contrario,
l'espansione topologica richiederebbe di modificare la topologia del grafo e ricalcolare le
\textit{feature} di tutti i nodi di sincronizzazione interessati.

\subsubsection{Operatività con Telemetria Parziale}
\label{sec:operativita_parziale}

Il \textit{framework} proposto non necessita di copertura totale dello spazio metrico teoricamente
possibile. Qualora alcune metriche non fossero disponibili, ad esempio per impossibilità di
instrumentare un \textit{database legacy} o dell'assenza di \textit{distributed tracing} su
componenti \textit{third-party}, il sistema è in grado di adattare il perimetro di analisi
attraverso la \textit{feature selection} guidata dall'ipergrafo. La Fase~I della
\textit{pipeline} seleziona automaticamente il sottoinsieme di metriche effettivamente osservabili
e rilevanti per la proprietà monitorata. In assenza di una metrica ideale, il \textit{forecasting}
opera su \textit{indicatori proxy}; ad esempio, se la saturazione del \textit{connection pool} non
è tracciabile in modo diretto, la latenza delle \textit{query} può essere utilizzata come
indicatore indiretto della contesa sulla risorsa. Nonostante la granularità e la completezza della
telemetria siano importanti, in quanto influenzano la profondità dell'analisi, esse non invalidano
l'architettura complessiva. In scenari con telemetria ridotta, infatti, l'approccio ibrido
continua a funzionare attraverso l'\textit{identificazione topologica delle interferenze}
(Layer~1) e il \textit{monitoraggio della Path Adherence} (Fase~III). Questi meccanismi consentono
infatti di mantenere una visibilità strutturale sul sistema anche in assenza di un'analisi causale
completa.

% =========================================================
\section{Pipeline di Forecasting della Non-Compliance}
\label{sec:pipeline_forecasting}
% =========================================================

La \textit{pipeline} di \textit{forecasting} della \textit{non-compliance} costituisce il motore
del \textit{framework} e si occupa di trasformare il flusso continuo di dati di telemetria in
stime quantitative relative al tempo residuo prima di una violazione di \textit{compliance}
all'interno del sistema in analisi. La \textit{pipeline} è costituita da quattro fasi sequenziali.
Le prime tre fasi si occupano di generare segnali informativi parziali, mentre l'ultima si occupa
di produrre un \textit{alert} semanticamente completo. Questo \textit{output} finale non si limita
ad avvisare della violazione prevista, bensì identifica e descrive anche la causa originaria della
violazione e il \textit{compliance set} interessato. L'efficacia della \textit{pipeline} risiede
principalmente nel modo in cui i tre \textit{layer} del \textit{framework} interagiscono in ogni
fase della \textit{pipeline} stessa.

\begin{itemize}
  \item \textbf{Layer~1 (Ipergrafo):} interviene nella Fase~I guidando in modo deterministico la
        \textit{feature selection} attraverso l'identificazione delle metriche rilevanti per ogni
        proprietà. Nella Fase~II, invece, orienta lo \textit{screening} causale verso le coppie
        di metriche semanticamente significative; infine, interviene nella Fase~IV dove la logica
        di aggregazione è determinata direttamente dalla semantica delle proprietà, la quale è
        codificata nell'ipergrafo;
  \item \textbf{Layer~2 (\gls{ATG}):} rappresenta il substrato computazionale di tutte e quattro
        le fasi. Ogni previsione, test causale o rilevamento di anomalia viene infatti eseguito
        sulle serie temporali contenute nel grafo temporale;
  \item \textbf{Layer~3 (\textit{Mapping}):} è operativo nelle Fasi~I e~II, dove la funzione
        $\mathcal{M}$ seleziona le metriche e le coppie da analizzare. Opera inoltre nella
        Fase~IV, dove la stessa funzione determina quali previsioni locali contribuiscono alla
        previsione globale della proprietà. Nella Fase~III, il Layer~3 non interviene direttamente;
        le analisi operano infatti sull'insieme $M_{\Phi_i}$ già costruito in Fase~I e sul \gls{PBO}
        definito nella Sezione~\ref{sec:pbo}.
\end{itemize}

Le fondamenta delle prime tre fasi sono costituite da tecniche di analisi consolidate nel campo dei
grafi e delle serie temporali. In questo caso, la differenza metodologica sostanziale rispetto a un
approccio tradizionale, non guidato semanticamente, risiede nel ruolo dell'ipergrafo. Quest'ultimo
orienta in modo deterministico la selezione delle metriche, lo \textit{screening} causale e
l'aggregazione dei segnali di anomalia. Questo approccio permette al \textit{framework} di
ottimizzare in modo significativo lo spazio di ricerca rispetto a un approccio puramente
statistico; così facendo, infatti, il \textit{framework} evita l'esplorazione completa delle
coppie di metriche, concentrando l'analisi sulle metriche topologicamente rilevanti per ciascuna
proprietà.

% ---------------------------------------------------------
\subsection{Fase~I: Previsione Locale delle Metriche}
\label{sec:fase_1}
% ---------------------------------------------------------

La prima fase della \textit{pipeline} ha il compito di trasformare le serie temporali di telemetria
grezza in previsioni quantitative riguardanti l'evoluzione delle metriche su un orizzonte temporale
$\tau$. L'output di questo stadio è rappresentato dall'insieme delle previsioni:
\begin{equation}
\label{eq:fase1_output}
\bigl\{\hat{X}^m(t+\tau') : m \in M_{\Phi_i},\; \tau' \in \{1,\ldots,\tau\}\bigr\}
\end{equation}
calcolate per ogni proprietà $\Phi_i$ e per ogni metrica appartenente al \textit{set} di
\textit{feature} selezionate.
Il parametro $\tau$ è pensato per essere configurabile in base alle necessità operative del sistema.
Attraverso un orizzonte temporale ridotto, ad esempio nell'ordine delle ore, è possibile garantire
una precisione superiore; tuttavia ciò risulta meno efficace per la pianificazione di interventi
preventivi complessi. Al contrario, orizzonti più estesi, nell'ordine di giorni o settimane,
producono preavvisi operativamente significativi ma con incertezza maggiore. Difatti, l'accuratezza
dei modelli di \textit{forecasting} degrada sistematicamente all'aumentare di $\tau$. Tuttavia, il
\textit{framework} è in grado di determinare automaticamente il \textit{lead time} $\tau^*$
individuando l'orizzonte minimo in cui la traiettoria prevista interseca la soglia certificata.
Questo meccanismo consente all'operatore di ricevere il preavviso massimo disponibile con una
determinata soglia di affidabilità della previsione.

\subsubsection{Feature Selection Topologica}
\label{sec:feature_selection}

In un sistema basato esclusivamente su \gls{ATG}, privo quindi del \textit{layer} semantico, la
\textit{feature selection} verrebbe trattata come un problema puramente statistico. Il numero di
coppie di metriche da analizzare crescerebbe in modo quadratico rispetto al numero totale di
metriche del sistema, raggiungendo per sistemi di medie dimensioni l'ordine di milioni di coppie.
Questo calcolo non solo richiederebbe \textit{dataset} storici estremamente ricchi, ma sarebbe
caratterizzato dall'incapacità intrinseca di rilevare interferenze \textit{cross-property} non
ancora manifestatesi nei dati.
Nell'approccio ibrido presentato, invece, la \textit{feature selection} è deterministica e guidata
dalla topologia del Layer~1. L'insieme $M_{\Phi_i}$ viene costruito integrando tre componenti
distinti.

\begin{itemize}
  \item \textbf{Componente Diretto ($M^{\text{direct}}$)}: include tutte le metriche di nodo e di
        arco associate ai componenti del \textit{compliance set} $H_{\Phi_i}$ e alle loro
        interazioni identificate dalla funzione di \textit{mapping}. Queste metriche sono
        considerate rilevanti per definizione;
  \item \textbf{Componente di Interferenza ($M^{\text{interf}}$)}: estende l'insieme del primo
        componente con le metriche di \textit{throughput} degli archi che convogliano traffico
        verso i nodi condivisi da altri \textit{compliance sets}. Tale traffico può infatti
        saturare le risorse condivise e impattare la proprietà $\Phi_i$, pur provenendo da un
        flusso non certificato per quest'ultima;
  \item \textbf{Componente Causale ($M^{\text{causal}}$)}: aggiunge, attraverso un'analisi causale
        preliminare \textit{offline}, le metriche di nodi esterni a $H_{\Phi_i}$ che mostrano
        dipendenze causali storicamente verificate con metriche interne al \textit{compliance set}.
        Affinché tali metriche possano essere scoperte e incluse nell'analisi, è requisito necessario
        che i nodi esterni siano instrumentati e appartengano all'infrastruttura di osservabilità
        globale del sistema, ossia l'insieme $V$ dell'\gls{ATG}. L'analisi non estende
        automaticamente il perimetro di strumentazione; difatti, se un nodo non è instrumentato,
        la sua dipendenza causale rimane un \textit{blind spot} noto del \textit{framework}.
\end{itemize}

L'insieme complessivo delle metriche da sottoporre a previsione è espresso formalmente come:
\begin{equation}
\label{eq:m_total}
M_{\Phi_i} = M_{\Phi_i}^{\text{direct}} \cup M_{\Phi_i}^{\text{interf}} \cup M_{\Phi_i}^{\text{causal}}
\end{equation}
Considerando la proprietà $\Phi_{\text{crit}}$, associata al percorso clinico
SGW$\to$SIN$\to$SDB$\to$SAL e agli archi $e_1$, $e_2$, $e_3$, le prime due componenti della
\textit{feature selection} vengono identificate dalla topologia. L'ipergrafo consente infatti di
aggirare lo \textit{screening} statistico, estraendo a priori le metriche rilevanti. Il componente
diretto include il vettore completo delle metriche per i tre archi del percorso e le metriche di
stato interno di tutti i quattro componenti appartenenti a $H_{\text{crit}}$. La selezione delle
metriche di nodo rispecchia le caratteristiche architetturali di ciascun servizio: per SGW e SAL,
che non gestiscono \textit{pool} di connessioni, le grandezze più informative sono l'utilizzo della
\gls{CPU} e l'occupazione della memoria; per SIN, che esegue operazioni di normalizzazione dei dati
biomedici su \gls{JVM}, si aggiungono le pause del \textit{garbage collector} come indicatore di
pressione sull'\textit{heap}; per SDB, che costituisce la risorsa condivisa critica, oltre al
consumo di \gls{CPU} e memoria viene monitorato il tasso di saturazione del \textit{connection
pool}. Complessivamente l'insieme risulta:
\begin{equation}
\label{eq:m_direct_crit}
M^{\text{direct}}(H_{\text{crit}}) =
  \bigcup_{e \in A(H_{\text{crit}})} \{L_e,\, \varepsilon_e,\, \theta_e\}
  \cup \bigcup_{v \in H_{\text{crit}}} \{\text{cpu}_v,\, \text{mem}_v\}
  \cup \{\text{gc}_{\text{SIN}},\, \text{pool}_{\text{SDB}}\}
\end{equation}
Il componente di interferenza include invece il \textit{throughput} dell'arco $e_5$
(SBL$\to$SDB), appartenente al flusso amministrativo: $M^{\text{interf}}(H_{\text{crit}}) =
\{\theta_{e_5}\}$.
L'arco $e_5$ viene incluso in quanto SDB è un nodo condiviso: il traffico proveniente da
$\Phi_{\text{adm}}$ può saturare il suo \textit{connection pool}, degradando il servizio verso SIN.
Senza il contributo semantico dell'ipergrafo, l'inclusione di questo arco richiederebbe la verifica
statistica della correlazione tra $e_5$ e le metriche di $H_{\text{crit}}$ sui dati storici. Se
tale interferenza non si fosse ancora verificata, un sistema \gls{ATG} puro ignorerebbe questo
parametro, lasciando il sistema cieco di fronte a una saturazione del \textit{connection pool}
causata da un improvviso picco di traffico nel flusso di fatturazione.

Il componente causale $M_{H_{\text{crit}}}^{\text{causal}}$, infine, emerge dall'analisi di Granger
eseguita sui dati storici per catturare dipendenze non esplicitate dalla topologia. Un esempio
significativo è il servizio SBL: pur non appartenendo al flusso clinico ($H_{\Phi_{\text{crit}}}$),
la sua attività di fatturazione può innescare contesa sulle risorse di I/O del disco condiviso con
SDB. Il test di Granger, applicato alle serie storiche di $\text{cpu}_{\text{SBL}}$ e di $L_{e_2}$,
potrebbe rilevare una dipendenza causale: periodi di elevata attività di SBL precedono
sistematicamente picchi di latenza sul flusso clinico. Questa dipendenza emerge dall'analisi
statistica e non dalla topologia dell'ipergrafo, giustificando la sua inclusione in
$M^{\text{causal}}$ anziché in $M^{\text{interf}}$. 
Questa dipendenza causale \emph{tra} $\text{cpu}_{\text{SBL}}$ e le metriche
del flusso clinico è distinta da quella già catturata topologicamente in
$M^{\text{interf}}$ tramite $\theta_{e_5}$: il Layer~1 intercetta
l'interferenza di carico sul \textit{connection pool}; il test di Granger
intercetta la contesa sulle risorse di I/O del disco, non direttamente
visibile nella topologia.

L'analisi che produce $M^{\text{causal}}$ è un'operazione eseguita in fase di \textit{setup} del
\textit{framework} sui dati storici disponibili e aggiornata periodicamente in seguito a modifiche
architetturali significative o in seguito alla rilevazione di derive comportamentali persistenti non
spiegabili dalle catene causali note. È concettualmente distinta dall'analisi causale della
Fase~II; difatti, la Fase~II opera a \textit{runtime} sull'insieme completo $M_{\Phi_i}$ che
include sia le metriche già note da $M^{\text{causal}}$ sia quelle identificate topologicamente;
in questo modo è quindi in grado di scoprire nuove relazioni causali emergenti dal comportamento
corrente del sistema, non solo validare quelle pre-identificate \textit{offline}.

\subsubsection{Modelli Predittivi Adottati}
\label{sec:modelli_predittivi}

Per ogni metrica $m \in M_{\Phi_i}$, il sistema applica un modello predittivo scelto in base alle
caratteristiche dinamiche della relativa serie temporale. Il \textit{framework} adotta una strategia
di \textit{forecasting} articolata su quattro livelli di complessità crescente, bilanciando
accuratezza e interpretabilità.

\paragraph{Regressione lineare come baseline}

Questo rappresenta l'approccio più semplice ed assume che la metrica segua un \textit{trend}
lineare nel tempo. La previsione consiste nella proiezione della retta di regressione, stimata
sulla finestra storica, sull'orizzonte $\tau$. La regressione lineare, basata sul metodo dei 
Minimi Quadrati Ordinari (OLS), rappresenta un riferimento eccellente in termini di interpretabilità 
e velocità; fornisce infatti una previsione che si possa facilmente verificare a livello visivo sul grafico 
della metrica. La rigidità è però il suo limite principale: non è in grado
di modellare la stagionalità, gestire i \textit{changepoint}\footnote{Nell'analisi delle serie temporali, 
un \textit{changepoint} indica una variazione improvvisa e discontinua del regime statistico dei dati, 
come un salto netto della latenza media a seguito di un deploy architetturale. 
Si distingue dalle dinamiche non lineari, che rappresentano invece evoluzioni continue ma irregolari, 
come la crescita esponenziale tipica di un \textit{memory leak}.} o catturare dinamiche non lineari.
Per questa ragione non viene adottata come modello principale, bensì come \textit{baseline} di
validazione. Qualora la divergenza tra la previsione del modello principale e quella della
\textit{baseline} superi una soglia configurabile, il sistema segnala un'incertezza elevata e
riduce la classificazione di criticità dell'eventuale \textit{alert} generato.

\paragraph{ARIMA per serie con memoria stocastica}

Il modello \gls{ARIMA}(p,d,q)~\cite{box2016} 
rappresenta un'evoluzione significativa rispetto alla regressione lineare. Difatti, mentre la
regressione assume che il futuro dipenda esclusivamente dal tempo, il modello ARIMA assume che il
valore futuro dipenda dalla storia recente della metrica stessa; questo avviene catturando la
memoria stocastica della serie. Il modello ARIMA è composto da tre diverse componenti. La
\textbf{componente autoregressiva (AR)} di ordine $p$ esprime il valore corrente come combinazione
lineare dei $p$ valori precedenti; la \textbf{componente di integrazione (I)} di ordine $d$ applica
$d$ differenziazioni successive per rendere stazionaria una serie che presenta un \textit{trend};
infine, la \textbf{componente a media mobile (MA)} di ordine $q$ incorpora gli errori di previsione
dei $q$ passi temporali precedenti. L'applicazione del modello segue quattro passi fondamentali.
Inizialmente avviene la verifica della stazionarietà attraverso il test di \gls{ADF} 
\footnote{Il test di Augmented Dickey-Fuller è un test statistico
standard utilizzato nell'analisi delle serie storiche per verificare l'ipotesi nulla di presenza
di una radice unitaria nel campione. Il superamento del test (con $p$-value tipicamente $< 0.05$)
permette di rigettare l'ipotesi nulla, confermando la stazionarietà della serie.}, con una soglia
$p$-value tipicamente pari a $0{,}05$. A questa prima fase segue una fase di differenziazione della
serie, se necessario, con ordine $d = 1$ o $d = 2$. Dopodichè avviene l'identificazione degli
ordini $p$ e $q$ attraverso i grafici di autocorrelazione (\gls{ACF}) e autocorrelazione parziale
(\gls{PACF}); infine, avviene la stima dei parametri per \gls{MLE}\footnote{La stima di \gls{MLE} è un metodo statistico classico utilizzato 
per inferire i parametri di un modello. Nel contesto del \textit{framework}, i \textit{solver} matematici 
sottostanti ai modelli predittivi (come ARIMA) impiegano la \gls{MLE} per calcolare i coefficienti ottimali, 
ovvero quelli che massimizzano la probabilità congiunta di osservare l'esatta sequenza dei dati storici di addestramento.}
e la validazione dei residui con il test di
Ljung-Box\footnote{Il test di Ljung-Box verifica l'assenza di autocorrelazione 
seriale nei residui: un risultato non significativo ($p > 0.05$) indica 
che i residui si comportano come rumore bianco, validando la bontà 
dell'adattamento del modello.}.

Il modello ARIMA è preferibile rispetto alla regressione per metriche che mostrano \textit{trend}
ben definiti e una componente stocastica significativa, soprattutto in presenza di serie temporali
sufficientemente lunghe, assenza di stagionalità complessa e necessità di intervalli di confidenza
solidi dal punto di vista statistico. Il limite principale risiede nella linearità; ARIMA non
cattura infatti dipendenze non lineari tra valori passati e futuri e gestisce in maniera poco
efficace \textit{changepoint} causati da eventi o anomalie di sistema.

\paragraph{Prophet per serie con stagionalità}

Molte metriche architetturali presentano \textit{pattern} fortemente ciclici legati ai ritmi di
utilizzo del sistema, come l'aumento del traffico durante le ore diurne. Per modellare queste serie
temporali caratterizzate da stagionalità multiple, il \textit{framework} utilizza l'algoritmo
Prophet~\cite{taylor2018forecasting}. A differenza dei classici modelli autoregressivi, Prophet
scompone il segnale secondo un approccio additivo basato su tre componenti principali. La prima
componente è rappresentata da un \textbf{\textit{trend} a lungo termine}, ed è in grado di
adattarsi ai cambiamenti strutturali del carico rilevando in modo del tutto autonomo i
\textit{changepoint}. Questa è modellata attraverso una funzione lineare o logistica, a seconda delle
caratteristiche della serie. La seconda è la \textbf{componente stagionale}, con periodicità
configurabile, la quale modella le fluttuazioni periodiche attraverso l'utilizzo di serie di Fourier.
La terza componente permette, infine, di integrare la conoscenza di eventi noti come
\textit{backup} notturni o rilasci di sistema, che altrimenti potrebbero inquinare temporaneamente
il comportamento della metrica.
La capacità di identificare autonomamente i \textit{changepoint} nel \textit{trend} fa sì che
Prophet risulti molto più robusto di ARIMA in scenari operativi dove il sistema può subire
interventi manuali. È inoltre in grado di gestire autonomamente i dati mancanti senza richiedere
interpolazione preventiva, poiché l'approccio di \textit{curve-fitting} non impone osservazioni su
una griglia temporale regolare. Per questi motivi, rappresenta la scelta più adatta per metriche
con una forte componente stagionale o influenzate da routine amministrative ricorrenti.

\paragraph{LSTM per dinamiche non lineari}

Per le metriche caratterizzate da dinamiche non lineari, come l'occupazione progressiva
dell'\textit{heap} \gls{JVM} in presenza di un \textit{memory leak} o la saturazione esponenziale
del \textit{connection pool}, il sistema ricorre a una rete \textbf{LSTM (\textit{Long
Short-Term Memory})}~\cite{hochreiter1997} addestrata sui dati storici della metrica in analisi.
Queste reti superano i limiti delle reti neurali ricorrenti tradizionali grazie a un meccanismo di
memoria selettiva, detto \textit{gate mechanism}, che determina in modo adattivo quali informazioni
della sequenza di \textit{input} debbano essere mantenute nello stato interno e quali possano
essere scartate in quanto irrilevanti per la previsione futura.
L'architettura prevede una struttura \textit{encoder-decoder}~\cite{sutskever2014}, dove l'\textit{encoder}
si occupa di comprimere la sequenza storica in uno stato interno vettoriale, mentre il
\textit{decoder} utilizza quello stato per generare le previsioni future sull'orizzonte temporale
$\tau$. La rete viene addestrata minimizzando una funzione di \textit{loss} appropriata al
\textit{task} di regressione; nel contesto del \textit{forecasting} di metriche continue, la
scelta tipica è l'Errore Quadratico Medio (\gls{MSE}), che penalizza in modo quadratico le
deviazioni sull'intero orizzonte predittivo.
Il \textit{framework} riserva la rete LSTM alle metriche con comportamenti fortemente non lineari o
dipendenze a lungo termine. L'elevato costo computazionale e la minore interpretabilità, comune ai
modelli di \textit{deep learning} rispetto alle controparti statistiche, la rendono infatti poco
adatta a un impiego generalizzato. Il significativo miglioramento in termini di accuratezza
predittiva, tuttavia, giustifica il suo utilizzo in presenza di dinamiche complesse o non lineari.
L'incertezza della previsione LSTM viene stimata tramite la tecnica \textit{Monte Carlo
Dropout}~\cite{gal2016}: durante l'\textit{inference}, il \textit{dropout} rimane attivo e la rete
viene eseguita $N$ volte sulla stessa sequenza di \textit{input}, producendo una distribuzione di
previsioni da cui si estraggono media e deviazione standard come stima dell'intervallo di
confidenza.

\begin{table}[htbp]
\centering
\caption{Confronto tra i modelli predittivi adottati nel framework.}
\label{tab:modelli_predittivi}
\small
\setlength{\tabcolsep}{4pt}
\makebox[\textwidth][c]{%
\begin{tabular}{@{}lcccc@{}}
\toprule
\textbf{Modello} & \textbf{Accuratezza} & \textbf{Interpretabilità} & \textbf{Costo comp.} & \textbf{Caso d'uso principale} \\
\midrule
Regressione lineare & Bassa     & Alta  & Trascurabile & \textit{Baseline}, validazione \\
ARIMA               & Media     & Media & Basso        & \textit{Trend} con memoria stocastica \\
Prophet             & Media-Alta & Media & Medio        & Stagionalità, dati irregolari \\
LSTM                & Alta       & Bassa & Alto         & Dinamiche non lineari \\
\bottomrule
% \multicolumn{5}{p{0.92\textwidth}}{\footnotesize\textit{Le valutazioni di accuratezza sono riferite al contesto applicativo del \textit{framework}, ossia serie temporali di metriche di sistemi distribuiti caratterizzate da componenti stagionali, \textit{trend} potenzialmente non lineari e dinamiche di saturazione. In contesti diversi, la gerarchia relativa tra modelli può variare.}} \\
\end{tabular}%
}
\end{table}

\subsubsection{Output della Fase~I}
\label{sec:output_fase1}

L'output della Fase~I consiste, per ogni proprietà $\Phi_i$ monitorata, nell'insieme delle
previsioni delle metriche appartenenti a $M_{\Phi_i}$ sull'orizzonte $\tau$. Per ogni serie
temporale selezionata, il sistema produce una traiettoria di previsioni sull'intero orizzonte da 1
a $\tau$, ciascuna accompagnata da un intervallo di confidenza che riflette l'incertezza del
modello. Ad esempio, in uno scenario in cui il flusso amministrativo eserciti una pressione
crescente sul \textit{database} SDB, l'output relativo alla proprietà $\Phi_{\text{crit}}$, su un
orizzonte ipotetico di 12 giorni, potrebbe produrre una previsione di latenza per l'arco critico
$e_2$ (SIN$\to$SDB) significativamente superiore alla soglia locale di 250\,ms adottata nelle
Tabelle~\ref{tab:scenario_memory_leak} e~\ref{tab:scenario_interferenza} per i singoli archi
critici. Le latenze degli altri archi del percorso, invece, risulterebbero ancora entro i limiti.

% ---------------------------------------------------------
\subsection{Fase~II: Identificazione delle Relazioni Causali}
\label{sec:fase_2}
% ---------------------------------------------------------

La seconda fase della \textit{pipeline} ha il compito di identificare le dipendenze causali tra le
metriche dell'insieme $M_{\Phi_i}$, distinguendo le correlazioni spurie dai rapporti di causa-effetto
verificati statisticamente. L'output di questo stadio è il \textbf{grafo causale}, ossia una
struttura orientata i cui nodi rappresentano le metriche analizzate, sia di nodo che di arco, i cui
archi rappresentano le dipendenze causali verificate. Ogni arco è etichettato con la tipologia di
dipendenza, lineare o non lineare, e l'intensità della relazione misurata. Per le relazioni lineari,
l'intensità è quantificata dall'incremento del coefficiente di determinazione $R^2$, del test di
Granger, tra il modello ristretto e il modello completo; per quelle non lineari, dal valore
\textit{TE} normalizzato. La distinzione tra correlazione e causalità è operativamente critica.
Due metriche possono essere fortemente correlate per diverse ragioni: la prima causa la seconda, la
seconda causa la prima, o perché entrambe sono l'effetto di una terza variabile nascosta. Un
operatore che riceve un \textit{alert} contenente una causa radice errata rischia di intervenire
sul componente sbagliato, finendo per consumare risorse senza risolvere il problema e lasciando
intatta la causa reale fino alla violazione effettiva dello \gls{SLA}.

\subsubsection{Screening Guidato dalla Topologia}
\label{sec:screening_topologia}

Analogamente a quanto avviene nella Fase~I, la struttura dell'ipergrafo permette di ridurre
drasticamente lo spazio di ricerca dell'analisi causale. Difatti, invece di calcolare le
correlazioni su tutte le coppie possibili di metriche presenti nel sistema, l'analisi si concentra
su tre categorie di coppie semanticamente rilevanti:

\begin{itemize}
  \item \textbf{Coppie intra-hyperedge}: riguardano le metriche su nodi e archi appartenenti allo
        stesso \textit{compliance set} $H_{\Phi_i}$. Queste metriche rappresentano le candidate
        naturali per l'identificazione di relazioni causali dirette tra lo stato interno di un nodo
        e la degradazione degli archi che esso serve;
  \item \textbf{Coppie inter-hyperedge}: coinvolgono metriche su nodi condivisi tra diversi \textit{compliance sets}
        $v \in \text{Shared}(H_{\Phi_i}, H_{\Phi_j})$ e metriche di archi appartenenti ad
        $A(H_{\Phi_j})$. Questa categoria è fondamentale per l'identificazione delle relazioni
        causali \textit{cross-property}, ossia come il carico di una proprietà possa influenzare
        le prestazioni di un'altra;
  \item \textbf{Coppie nodo-arco}: ossia le coppie dove il nodo $v$ appartiene a $H_{\Phi_i}$ e
        l'arco $e$ appartiene ad $A(H_{\Phi_i})$. Le coppie nodo-arco costituiscono il
        sottoinsieme più rilevante delle coppie intra-hyperedge ai fini della \textit{root cause
        analysis}, in quanto collegano direttamente lo stato interno di un componente alla qualità
        dell'interazione che esso serve.
\end{itemize}

Le tre categorie non sono mutualmente esclusive per costruzione; la distinzione serve a identificare
le priorità dell'analisi.

\subsubsection{Procedura di Analisi Causale in Tre Stadi}
\label{sec:analisi_causale}

Per ogni coppia di metriche appartenente alle categorie descritte, l'analisi causale si articola in
tre fasi sequenziali di complessità crescente, bilanciando rigore statistico ed efficienza
computazionale.

\paragraph{Screening preliminare (Correlazione di Pearson)}

Il primo stadio consiste in un filtraggio basato sul coefficiente di correlazione lineare di
Pearson $r \in [-1, 1]$. Invece di eseguire immediatamente test causali complessi, il \textit{framework}
promuove al test di Granger della fase successiva solo le coppie che mostrano un legame
relazionale forte, applicando una soglia operativa\footnote{La soglia operativa $|r| > 0.7$ 
funge da euristica per il pre-filtraggio: permette di scartare relazioni deboli o dominate dal 
rumore prima di impegnare risorse computazionali nei più costosi test di causalità temporale.}, 
ad esempio $|r| > 0{,}7$. Questa scelta
introduce un \textit{trade-off} tra efficienza computazionale e rischio di falsi negativi. Il
filtraggio permette infatti di ridurre drasticamente il numero di test successivi, aumentando di
conseguenza l'efficienza computazionale. D'altra parte, dipendenze non lineari come gli effetti
soglia, tipici dei sistemi distribuiti, possono non manifestarsi attraverso una correlazione lineare
significativa: tali relazioni presentano spesso valori di correlazione lineare di Pearson bassi
prima della soglia e crescono rapidamente solo dopo il superamento di essa, risultando quindi non
idonee a passare allo stadio successivo. Per mitigare questo rischio, il sistema introduce un
meccanismo di \textit{bypass} guidato dalla topologia. Le coppie di metriche identificate dal
Layer~1 come \textit{inter-hyperedge}, ovvero associate a nodi condivisi tra diversi \textit{compliance
sets}, vengono sottoposte direttamente al test di Granger \textit{bypassando} il filtro di Pearson.
Questa scelta è motivata dal fatto che, per un insieme di coppie già ridotto per costruzione, il
filtro di Pearson non produce un risparmio computazionale significativo, mentre introduce il rischio
concreto di scartare relazioni causali non lineari prima che possano essere valutate dagli stadi
successivi.

\paragraph{Test di causalità di Granger}

Le coppie che superano il primo filtro vengono successivamente sottoposte all'analisi di causalità
di Granger~\cite{granger1969}. Questo approccio formalizza un principio predittivo: se una metrica
$m_1$ causa effettivamente una metrica $m_2$, allora la conoscenza della storia passata di $m_1$
deve migliorare la precisione nel prevedere il futuro di $m_2$ rispetto all'uso della sola storia
di $m_2$. A livello operativo, il sistema confronta un modello ``ristretto'' con un modello
``completo''. Il modello ristretto si basa esclusivamente sui valori passati della metrica in analisi
mentre, il modello completo, integra i valori passati della potenziale causa. Se il miglioramento
apportato dal modello completo rispetto a quello ristretto risulta statisticamente significativo,
attraverso un test $F$ sulla riduzione della somma dei quadrati dei residui tra il modello completo
e quello ristretto con $p$-value $< 0{,}05$, è possibile concludere che $m_1$ \textbf{causa-Granger}
$m_2$, ovvero che la sua storia passata ha valore predittivo su $m_2$.
Il test richiede che le serie temporali analizzate siano stazionarie. Per le metriche per cui
l'analisi in Fase~I ha determinato un ordine di differenziazione $d > 0$ tramite il test \gls{ADF},
la medesima trasformazione viene applicata alla serie storica prima del test di Granger. Per le
metriche trattate con modelli diversi da ARIMA, il test \gls{ADF} viene eseguito preventivamente
con la stessa procedura descritta in Fase~I.

\paragraph{Transfer Entropy Non-lineare}

Il terzo stadio funge da meccanismo di \textit{fallback} per quelle coppie che, pur avendo
superato lo \textit{screening} iniziale di Pearson o essendone state esentate, come le coppie
\textit{inter-hyperedge}, non superano il test di Granger, tale per cui la loro dipendenza causale
potrebbe non essere lineare. In questi casi, il \textit{framework} applica la \textit{Transfer
Entropy} ($\text{TE}_{m_1 \to m_2}$)~\cite{schreiber2000}, ossia una metrica basata sulla teoria
dell'informazione che quantifica quanto la conoscenza di $m_1$ riduca l'incertezza su $m_2$. A
differenza del test di Granger, strettamente legato a modelli lineari, la \textit{Transfer Entropy}
estende l'analisi alle relazioni non lineari. Tale caratteristica è fondamentale, oltre che per
intercettare relazioni non lineari, per individuare i tipici ``effetti soglia'' dei microservizi,
ossia situazioni in cui le prestazioni crollano improvvisamente una volta raggiunta una certa
saturazione. La messa in atto di questa terza fase è selettiva; ciò è dovuto all'elevato costo
computazionale e alla necessità di discretizzare le serie temporali nell'implementazione basata
sullo stimatore classico.

\subsubsection{Identificazione della Causalità Cross-Property}
\label{sec:causalita_cross_property}

Nel caso di nodi condivisi tra diversi \textit{compliance sets}, il sistema verifica l'esistenza di
una catena causale \textit{cross-property} capace di collegare il carico di una proprietà alla
degradazione di un'altra. La forma di questa catena è specifica: un arco appartenente al
\textit{compliance set} di $\Phi_j$ causa la saturazione del nodo condiviso $v$, che a sua volta
causa la degradazione di un arco del \textit{compliance set} di $\Phi_i$. La verifica avviene
applicando test di Granger in sequenza lungo le due direzioni della catena; qualora entrambi
restituiscano esito positivo, l'interferenza \textit{cross-property} viene confermata statistica\-mente
e integrata nel grafo causale e nell'\textit{alert} della Fase~IV. \\
Questa procedura costituisce una verifica sequenziale della catena causale, non una Granger
causality condizionale in senso formale. Non esclude la presenza di variabili confondenti o percorsi
diretti alternativi. Viene adottata come compromesso tra rigore statistico e praticabilità
computazionale in contesti operativi, e costituisce evidenza consistente di dipendenza causale, non
una prova formale di causalità strutturale.

Nello scenario di interferenza del sistema di telemedicina, i test potrebbero rivelare una forte
correlazione tra il \textit{throughput} del flusso amministrativo e la latenza del flusso clinico.
In particolare, si potrebbe osservare una causalità di Granger dal \textit{throughput} su $e_5$
verso la saturazione del \textit{pool} di SDB, e una seconda causalità di Granger dalla saturazione
del \textit{pool} di SDB verso la latenza su $e_2$. La catena causale risultante,
$\theta_{e_5} \to \text{pool}_{\text{SDB}} \to L_{e_2}$, attraverserebbe due \textit{compliance
sets} distinti. Un sistema di monitoraggio tradizionale, operando su ciascuna proprietà in
isolamento, rimarrebbe cieco di fronte a questa connessione. Sarebbe in grado di identificare il
sintomo, ossia l'aumento della latenza clinica, ma non saprebbe risalire alla causa radice in
quanto quest'ultima risiede al di fuori del perimetro clinico monitorato.

% ---------------------------------------------------------
\subsection{Fase~III: Monitoraggio dell'Aderenza Strutturale}
\label{sec:fase_3}
% ---------------------------------------------------------

La terza fase della \textit{pipeline} rileva anomalie nello stato aggregato del \textit{compliance
set} e deviazioni strutturali del comportamento del sistema rispetto al percorso certificato.
L'analisi si articola su due tracce parallele che convergono in un validatore finale. La prima
traccia comprende il monitoraggio metrico puntuale ed è articolata su due livelli:
il livello base, composto da \textit{threshold} statico e z-score sempre attivi, e l'Isolation Forest che viene
attivato solo in presenza di almeno uno di questi segnali di base. La seconda traccia monitora costantemente la deviazione
comportamentale tramite \gls{CUSUM} sul \textit{Path Adherence Score}, in parallelo e in modo
indipendente dalla prima. Il livello di validazione strutturale finale viene attivato solo quando
entrambe le tracce segnalano un'anomalia. Il carico computazionale viene così ottimizzato senza
compromettere la sensibilità del sistema.

\subsubsection{Livello Base: Threshold Statico e Z-Score}
\label{sec:threshold_zscore}

Il primo livello di analisi, costantemente attivo, funge da ``sentinella'' attraverso l'applicazione
di un filtro deterministico sulle metriche contenute in $M_{\Phi_i}$. Questa fase combina la
rigidità necessaria dei vincoli contrattuali con la flessibilità della statistica di base.

\paragraph{Threshold statico: il vincolo contrattuale}

Per ogni metrica $m$ e ogni soglia $T_m$ definita nel modello di certificazione, viene
generato un \textit{alert} immediato al verificarsi di specifiche condizioni. Nel dettaglio, 
l'allarme scatta quando $m(t) > T_m$ per le metriche con limite superiore, come latenza o
\textit{error rate}, oppure quando $m(t) < T_m$ per metriche con limite inferiore, come affidabilità
e \textit{throughput}. Questo metodo traduce i requisiti di \textit{compliance} in un semplice 
confronto aritmetico, non richiede alcun addestramento ed è efficace nel rilevare violazioni nette. 
Tuttavia, il suo limite è dato dalla rigidità: non è in grado di distinguere tra un carico 
fisiologico previsto, come un picco durante le ore di punta, e un'anomalia strutturale che prelude a un guasto.

\paragraph{Z-score adattivo: il contesto statistico}

Per superare la rigidità del \textit{threshold} statico, il \textit{framework} adotta uno
\textbf{z-score adattivo}. Una metrica $m$ al tempo $t$ è segnalata come anomala se il suo valore
si discosta dalla propria \emph{media mobile storica} $\mu_m$, calcolata su una finestra
configurabile di osservazioni recenti, per un valore maggiore di un parametro configurabile di
deviazioni standard $\sigma_m$.
% \footnote{Lo z-score quantifica la distanza di una singola
% osservazione dalla propria media mobile storica in termini di deviazioni standard ($\sigma$). Sotto
% l'assunzione di una distribuzione normale, l'utilizzo di una soglia $|z| > 3$ garantisce un'elevata
% sensibilità, poiché isola unicamente i valori che ricadono in quel $0{,}3\%$ di casistiche estreme
% rispetto al comportamento nominale.}. 
Tipicamente, il sistema adotta una soglia $|z| > 3$ che, in
una distribuzione gaussiana, corrisponde ai valori che hanno una probabilità inferiore allo
$0{,}3\%$ di verificarsi in condizioni normali. Il vantaggio dello z-score è dato dalla
maggior sensibilità ai contesti operativi rispetto al \textit{threshold} statico; difatti, la soglia
dello z-score si calibra automaticamente sui dati storici, imparando di conseguenza il
comportamento ritenuto ``normale'' del sistema. Il suo limite risiede nell'assunzione implicita di
stazionarietà della distribuzione, che può non valere per metriche con \textit{trend} progressivi;
in questi casi, lo z-score fatica a distinguere tra anomalie reali e variazioni strutturali del
sistema. Per questa ragione viene utilizzato come segnale di primo livello per attivare le analisi
successive, non come strumento di diagnosi finale.

\subsubsection{Rilevamento di Anomalie Aggregate: Isolation Forest}
\label{sec:isolation_forest}

Il secondo livello applica l'Isolation Forest sullo stato aggregato del \textit{compliance set}
$H_{\Phi_i}$. Mentre i filtri del primo livello operano in modo isolato, l'\textbf{Isolation
Forest}~\cite{liu2008} analizza le correlazioni implicite tra i componenti del \textit{set} per
identificare anomalie multivariate.

\paragraph{Il vettore di stato aggregato}

Per ogni \textit{compliance set}, il sistema costruisce un vettore di stato aggregato concatenando
le \textit{feature} di nodo di tutti i componenti del \textit{set}:
\begin{equation}
\label{eq:aggregate_state}
\mathbf{X}_{H_{\Phi_i}}(t) = \bigl[\mathbf{x}_{v_1}(t),\; \mathbf{x}_{v_2}(t),\; \ldots,\; \mathbf{x}_{v_k}(t)\bigr]
\end{equation}
dove $k = |H_{\Phi_i}|$ rappresenta la cardinalità del \textit{set}. Su questo vettore aggregato
viene successivamente applicato un modello Isolation Forest, addestrato sui dati storici di periodi
di normale operatività. Le dimensioni del vettore corrispondenti a metriche strutturalmente non
rilevanti per un dato nodo, come $\text{pool}_v$ per SGW e SAL o $\text{gc}_v$ per SDB, assumono
valori costanti o prossimi allo zero; il modello Isolation Forest apprende questa distribuzione
durante la calibrazione sul \textit{gold standard}, incorporandola come parte del profilo di
normalità. Le \textit{feature} di arco vengono escluse dal vettore aggregato in quanto effimere:
in finestre temporali con traffico ridotto o assente, i loro valori sarebbero strutturalmente nulli
indipendentemente dallo stato di salute del sistema, rendendo instabili sia l'addestramento che
l'\textit{inference} del modello. L'Isolation Forest opera quindi esclusivamente sulle
\textit{feature} di nodo, che per costruzione assumono valori informativi anche in assenza di
traffico.

\paragraph{Principio di funzionamento: l'isolamento della rarità}

A differenza dei metodi tradizionali che cercano di definire cosa sia ``normale'', l'Isolation
Forest si concentra sull'identificazione diretta delle anomalie basandosi sulla \emph{rarità} e
sulla presenza di valori \emph{distanti} dalla distribuzione centrale. Il processo segue una
logica basata su alberi decisionali. Le \textit{osservazioni anomale}, essendo statisticamente rare,
si trovano in regioni dello spazio delle \textit{feature} poco popolate; ciò fa sì che possano
essere isolate in pochissimi passi. Le \textit{osservazioni normali}, invece, essendo dense e vicine
tra loro, richiedono numerose partizioni ricorsive prima di essere separate dalle altre. Una
\emph{profondità di isolamento bassa}, ossia un punto che si separa rapidamente dal resto della
foresta, è un segnale di atipicità e quindi di anomalia. Il sistema adotta una soglia di anomalia
$s > 0{,}5$, calibrabile tramite il parametro di \textit{contamination} disponibile nelle librerie
standard come scikit-learn\footnote{\url{https://scikit-learn.org}}, che definisce la frazione
attesa di anomalie nel \textit{set} di addestramento. Il vantaggio di applicare l'Isolation Forest
allo stato aggregato, anziché alle singole metriche, è che il modello è in grado di rilevare
combinazioni di valori che presi singolarmente apparirebbero conformi. Nell'ambito dello scenario
di riferimento, un \textit{database} SDB con \gls{CPU} al 15\% e \textit{pool} al 75\% potrebbe
apparire nei limiti qualora le due metriche venissero analizzate separatamente. Tuttavia, la loro
combinazione risulta fortemente anomala rispetto ai \textit{pattern} storici che, solitamente,
mostrano o uno stato normale avente entrambe le percentuali basse, o uno stato di picco con
entrambe le metriche elevate. Una combinazione di questo tipo segnala un'anomalia strutturale
profonda, come un \textit{connection leak} o un blocco I/O, che i filtri univariati ignorerebbero.

\subsubsection{Monitoraggio della Deriva Strutturale con EWMA e \gls{CUSUM}}
\label{sec:ewma_cusum}

Il terzo livello monitora la deviazione progressiva del comportamento del sistema rispetto al
percorso certificato, attraverso due strumenti complementari: \gls{EWMA} e \gls{CUSUM}.
L'obiettivo dell'analisi è raggiungere la capacità, da parte del \textit{framework}, di distinguere
gli scostamenti momentanei dei dati da un vero e proprio degrado strutturale.

\paragraph{Pulizia del segnale: il filtro \gls{EWMA}}

Prima di applicare il \gls{CUSUM}, il \textit{Path Adherence Score} viene sottoposto a un processo
di livellamento tramite il filtro \gls{EWMA}\footnote{\gls{EWMA} è un filtro statistico di tipo 
\gls{IIR} che applica un peso esponenzialmente decrescente alle osservazioni storiche, conferendo 
ai dati recenti maggiore rilevanza rispetto a quelli passati.} al fine di ridurre il rumore di fondo. Il segnale grezzo, infatti,
può oscillare significativamente tra finestre temporali consecutive, anche in assenza di anomalie
gravi in corso. \gls{EWMA} si basa su un fattore di smorzamento, con valore compreso tra 0 e 1,
che regola il peso attribuito ai dati passati permettendo di bilanciare tra reattività e stabilità
del segnale. In questo contesto, la reattività indica la capacità del sistema di segnalare
immediatamente un cambiamento; mentre la stabilità rappresenta la capacità di assorbire piccole
fluttuazioni casuali tra i dati, evitando di lanciare allarmi in maniera impropria.
Sostanzialmente, questo fattore determina la velocità con cui i dati più vecchi perdono importanza
rispetto a quelli nuovi: valori più alti fanno sì che i dati passati perdano importanza rapidamente,
rendendo il filtro più reattivo alle variazioni recenti; valori più bassi, invece, preservano la
memoria storica più a lungo, producendo una stima più stabile ma meno sensibile ai cambiamenti.
Calibrando correttamente il filtro è possibile ottenere una serie temporale omogenea e quindi adatta
a un monitoraggio continuo. Lo stesso trattamento di livellamento viene applicato alla serie
temporale della norma di Frobenius nei casi in cui il \gls{PAS} non risulti applicabile, garantendo
la medesima stabilità del segnale indipendentemente dalla metrica di aderenza utilizzata.

\paragraph{Rilevazione della deriva persistente: il metodo \gls{CUSUM}}

Sulla serie temporale filtrata attraverso \gls{EWMA} viene poi applicato l'algoritmo
\gls{CUSUM}~\cite{montgomery2009}, ossia uno strumento progettato per evidenziare derive strutturali
persistenti, distinguendole da fluttuazioni isolate. A differenza dei meccanismi di soglia
tradizionali, il \gls{CUSUM} non reagisce a singoli eventi anomali; bensì accumula progressivamente
le deviazioni negative del \textit{Path Adherence Score} rispetto al valore di riferimento
calcolato sul \textit{gold standard} $W^{\text{gold}}$. Al fine di ridurre il rischio di falsi
positivi, garantendo quindi solidità operativa, il meccanismo di accumulo si basa su un margine di
tolleranza configurabile. Questo parametro permette al sistema di assorbire le micro oscillazioni
fisiologiche del traffico, impedendo che variazioni irrilevanti alimentino il contatore
dell'anomalia. Il sistema genera un \textit{alert} di deviazione comportamentale non appena il
degrado accumulato supera una soglia predefinita, segnalando quindi un allontanamento significativo
e persistente dal percorso certificato.

L'impostazione del \gls{CUSUM} all'interno del \textit{framework} è volutamente unidirezionale,
ovvero orientata alla sola direzione di degrado della metrica su cui opera: per il \textit{Path
Adherence Score}, ciò significa accumulare i decrementi rispetto al valore di riferimento; per la
norma di Frobenius, significa accumulare gli incrementi rispetto al valore di riferimento calcolato
sul \textit{gold standard}. In entrambi i casi, il sistema si concentra esclusivamente
sull'allontanamento dal comportamento certificato; scostamenti nella direzione opposta, come un
picco di traffico che scorre correttamente lungo gli archi certificati, non costituiscono una
violazione della \textit{compliance} strutturale nel contesto qui definito.

\subsubsection{Verifica della Conformità Strutturale}
\label{sec:conformita_strutturale}

Il quarto livello, essendo il più oneroso dal punto di vista computazionale, viene attivato
esclusivamente come validatore finale. Il suo compito è confrontare il \gls{PBO} osservato con il
\textit{gold standard} attraverso una procedura articolata in tre fasi.

\begin{enumerate}
  \item \textbf{Analisi del \textit{trend}}: viene calcolata la derivata temporale della serie
        livellata tramite \gls{EWMA}, al fine di identificare la direzione del degrado. Per il
        \textit{Path Adherence Score}, il degrado si manifesta come decremento persistente, ossia
        come derivata negativa, mentre per le metriche di distanza strutturale, quali norma di
        Frobenius o \gls{GED}, il degrado si manifesta come incremento persistente, ossia derivata
        positiva.
  \item \textbf{Misurazione della distanza}: nel caso in cui la topologia sia fissa, ossia $V$ e
        $E_{\text{all}}$ invariati, come nello scenario di riferimento, la distanza strutturale
        viene elaborata attraverso la norma di Frobenius $\|W(t) - W^{\text{gold}}\|_F$, calcolata
        sull'intera matrice dei pesi. In contesti di topologia variabile a \textit{runtime},
        invece, la distanza viene calcolata attraverso la \textbf{\textit{Graph Edit Distance}
        (\gls{GED})} tra il \gls{PBO} corrente e il \textit{gold standard} per quantificare la
        deformazione del grafo.
  \item \textbf{Validazione dell'\textit{alert}}: qualora la derivata confermi per almeno tre
        intervalli consecutivi la direzione del degrado, derivata negativa per il \textit{Path
        Adherence Score} e positiva per la norma di Frobenius o \gls{GED}, e il valore della
        distanza strutturale superi una soglia configurabile, viene generato un segnale di
        deviazione strutturale confermata.
\end{enumerate}

La combinazione delle due condizioni che stanno alla base della terza fase, ossia la validazione
dell'\textit{alert}, è fondamentale per filtrare i falsi positivi causati da fluttuazioni temporanee
della metrica di aderenza comportamentale che non corrispondono a deformazioni strutturali effettive
del grafo comportamentale.

\subsubsection{Procedura di Rilevamento Combinata}
\label{sec:rilevamento_combinato}

I quattro livelli della Fase~III vengono combinati secondo un protocollo di attivazione gerarchico.
Il \textit{threshold} statico e lo z-score sono sempre in funzione e costituiscono il primo filtro
di selezione. L'Isolation Forest viene interrogato solo quando uno dei due segnali di base è attivo,
in modo tale da verificare se lo stato aggregato del \textit{compliance set} sia coerente con il
profilo storico. Parallelamente, il \gls{CUSUM} monitora costantemente la serie livellata tramite \gls{EWMA}: 
il \textit{Path Adherence Score} nei casi di percorso lineare ordinato, oppure la norma di Frobenius 
$\|W(t) - W^{\text{gold}}\|_F$ come segnale di deviazione scalare qualora il \gls{PAS} non risulti applicabile.
Un'allerta di \textit{drift} viene generata quando l'accumulo statistico di tali scostamenti supera 
la soglia critica. Infine, qualora sia il \gls{CUSUM} sia l'Isolation Forest segnalino anomalie, 
interviene il \textbf{validatore strutturale}. Quest'ultimo calcola puntualmente la distanza strutturale 
definitiva, tramite norma di Frobenius per topologie fisse o \gls{GED} per topologie variabili, 
per confermare formalmente la deviazione prima dell'emissione dell'alert.


La diagnosi finale dipende dalla combinazione dei segnali attivi. Se tutti e quattro gli indicatori sono
attivi simultaneamente, il sistema classifica lo stato come criticamente anomalo sia a livello
metrico che comportamentale; in questo caso la criticità dell'\textit{alert} in uscita dalla
Fase~IV viene elevata al livello massimo. Qualora solo il \gls{CUSUM} sia attivo, il sistema
segnala una deviazione comportamentale in corso che non ha ancora generato anomalie metriche
evidenti. Si tratta tipicamente di interferenze \textit{cross-property} nelle fasi iniziali.
Qualora invece il segnale di base sia attivo e il \gls{CUSUM} segnali una deviazione comportamentale, 
ma l'Isolation Forest non confermi l'anomalia multivariata, il sistema rileva la coesistenza di un'anomalia 
metrica puntuale e di una deriva comportamentale strutturale non caratterizzata da pattern aggregati anomali. 
In questo caso il Validatore non viene attivato; in assenza di incertezza modellistica elevata, 
la classificazione in Fase~IV sarà almeno Arancione per l'attività del \gls{CUSUM}, indipendentemente dal \textit{lead time}.
Altrimenti, qualora il \gls{CUSUM} e l'Isolation Forest segnalino entrambi un'anomalia, viene attivato
il validatore strutturale. Se quest'ultimo non conferma la deviazione, ad esempio in quanto la
distanza strutturale è ancora sotto soglia o la derivata non sia persistente, il sistema produce 
un segnale di preallerta dovuto al fatto che sono presenti anomalie sia metriche
che comportamentali, ma la deformazione strutturale del grafo non è ancora statisticamente confermata. 
La classificazione finale dell'\textit{alert} dipenderà poi dal \textit{lead time} 
stimato nella Fase~IV, che determinerà il livello Arancione o Rosso secondo le soglie definite nella 
Sezione~\ref{sec:classificazione_criticita}, in attesa della conferma strutturale 
del Validatore. 
% Altrimenti, qualora solo l'Isolation Forest sia attivo, il sistema segnala uno stato metrico anomalo
% che non ha ancora causato una deviazione di percorso. Qualora il segnale di base, \textit{threshold}
% o z-score, sia attivo ma l'Isolation Forest non confermi l'anomalia multivariata, il sistema
% registra un evento di attenzione a bassa priorità senza elevarlo ad \textit{alert} strutturale.
% Si tratta tipicamente di una fluttuazione di carico fisiologica, coerente con i \textit{pattern}
% storici: un picco locale di risorse che il sistema sta ancora gestendo attraverso il percorso
% certificato. Questa stratificazione consente di distinguere tra semplici fluttuazioni di carico e
% derive strutturali profonde, rendendo gli interventi più mirati e tempestivi.
Altrimenti, qualora l'Isolation Forest confermi l'anomalia multivariata e il segnale di base, 
\textit{threshold} o z-score, sia attivo, ma il \gls{CUSUM} non segnali deviazioni comportamentali, 
il sistema segnala uno stato metrico anomalo che non ha ancora causato una deviazione di percorso. 
Si tratta tipicamente di un picco di carico interno al \textit{compliance set} che non ha ancora 
alterato la distribuzione del traffico sul percorso certificato.
Infine, qualora il segnale di base sia attivo ma l'Isolation Forest non confermi l'anomalia multivariata, 
il sistema registra un evento di attenzione a bassa priorità senza elevarlo ad \textit{alert} strutturale.
Si tratta tipicamente di una fluttuazione di carico fisiologica, coerente con i \textit{pattern} storici.
Questa stratificazione consente di distinguere tra semplici fluttuazioni di carico e derive strutturali 
profonde, rendendo gli interventi più mirati e tempestivi.


% ---------------------------------------------------------
\subsection{Fase~IV: Generazione degli Alert e Stima del Lead Time}
\label{sec:fase_4}
% ---------------------------------------------------------

La quarta fase della \textit{pipeline} rappresenta il cuore del \textit{framework} proposto. Essa
è fondamentale in quanto aggrega i risultati delle tre fasi precedenti al fine di produrre
\textit{alert} relativi alla proprietà certificata. La differenza rispetto a un sistema di
\textit{anomaly detection} tradizionale è inizialmente quantitativa, in quanto permette di ottenere
un \textit{lead time} più ampio; tuttavia, il vantaggio principale è qualitativo: l'\textit{alert}
generato, infatti, rappresenta un vero e proprio giudizio di \textit{compliance} predittivo su una
specifica proprietà. Il \textit{framework}, attraverso la fusione delle previsioni locali (Fase~I),
delle catene causali (Fase~II) e delle deviazioni strutturali (Fase~III), permette infatti di
fornire un verdetto che include l'indicazione esplicita della causa strutturale e dei componenti
impattati.

\subsubsection{Aggregazione Predittiva per Proprietà}
\label{sec:aggregazione_predittiva}

Per ogni proprietà $\Phi_i$, la previsione globale viene calcolata applicando una funzione di
aggregazione semantica alle traiettorie predittive locali ottenute nella Fase~I:
\begin{equation}
\label{eq:global_forecast}
\hat{\Phi}_i(t+\tau') = f_{\text{agg}}^{\Phi_i}\!\bigl(\{\hat{X}^m(t+\tau') : m \in M_{\Phi_i}\}\bigr)
\end{equation}
La scelta della funzione $f_{\text{agg}}^{\Phi_i}$ dipende strettamente dalla natura della
proprietà monitorata. Il \textit{framework} copre le tre casistiche più comuni. Queste funzioni sono
definite per catene seriali, che costituiscono il caso architetturale del sistema di telemedicina
analizzato. In presenza di topologie con rami paralleli, la funzione di aggregazione
per la latenza diventerebbe ad esempio il massimo tra i percorsi paralleli anziché la somma.

Per semplificare la notazione nelle formule successive, si definisce $A_i = A(H_{\Phi_i})$ come 
l'insieme degli archi appartenenti al \textit{compliance set} della la proprietà $\Phi_i$, 
ovvero gli archi del grafo temporale i cui estremi appartengono entrambi a $H_{\Phi_i}$.
In questo formalismo, $t$ rappresenta l'istante temporale corrente; 
mentre $\tau$ indica l'orizzonte temporale massimo di previsione, ad esempio un intervallo di 12 giorni. 
La variabile $\tau' \in \{1, \ldots, \tau\}$ funge da indice per i singoli passi previsionali futuri. 
Tutte le quantità si intendono dunque valutate al tempo futuro $t+\tau'$; l'utilizzo di tale variabile 
è necessario per monitorare l'evoluzione dell'intera traiettoria predittiva e identificare con 
precisione l'istante in cui si prevede possa verificarsi una violazione, anziché limitare 
l'analisi al solo stato finale del sistema all'orizzonte $\tau$.

Per proprietà di latenza \textit{end-to-end}, la latenza totale è calcolata come la somma delle
latenze previste sui singoli archi del percorso certificato:
\begin{equation}
\label{eq:agg_latency}
\hat{L}_{\Phi_i}(t+\tau') = \sum_{e \in A_i} \hat{L}_e(t+\tau')
\end{equation}

Per proprietà di affidabilità, una richiesta \textit{end-to-end} ha successo solo se tutte le
interazioni della catena certificata non producono errori; questo riflette direttamente la semantica
AND del \textit{compliance set}. L'affidabilità globale prevista è il prodotto delle probabilità di
successo dei singoli archi:
\begin{equation}
\label{eq:agg_reliability}
\hat{R}_{\Phi_i}(t+\tau') = \prod_{e \in A_i} (1 - \hat{\varepsilon}_e(t+\tau'))
\end{equation}
dove $\hat{\varepsilon}_e(t+\tau')$ è il tasso di errore previsto sull'arco $e$, direttamente
forecastabile dalla componente $\varepsilon_e$ del vettore $\mathbf{x}_e(t)$ definito nella
Sezione~\ref{sec:def_formale_atg}.

Per proprietà di capacità, il limite della catena è determinato dal suo collo
di bottiglia, ossia dall'arco con \textit{throughput} minimo. La capacità globale prevista è:
\begin{equation}
\label{eq:agg_capacity}
\hat{C}_{\Phi_i}(t+\tau') = \min_{e \in A_i} \hat{\theta}_e(t+\tau')
\end{equation}
dove $\hat{\theta}_e(t+\tau')$ è il \textit{throughput} previsto sull'arco $e$, forecastabile
direttamente dalla componente $\theta_e$ del vettore $\mathbf{x}_e(t)$.

Nello scenario di riferimento, tra il flusso clinico e quello amministrativo, l'aggregazione
additiva della latenza su un orizzonte temporale di $\tau = 12$ giorni potrebbe indicare un valore
totale $\hat{L}_{\text{crit}}$ superiore alla soglia di sicurezza di 500\,ms. In questo
caso, il sistema procederebbe a segnalare preventivamente una violazione prevista, permettendo di
intervenire prima che il degrado diventi effettivo.

\subsubsection{Stima del Lead Time}
\label{sec:stima_lead_time}

Il \textit{lead time} $\tau^*$ associato a ciascun \textit{alert} indica il minimo
orizzonte temporale entro il quale la proprietà $\Phi_i$ monitorata violerà le condizioni di conformità.
Essendo che la violazione può corrispondere sia al superamento di un limite superiore, ad esempio la
latenza elevata, sia al cedimento al di sotto di un limite inferiore, come il calo dell'affidabilità
o del \textit{throughput}, la definizione di \textit{lead time} viene generalizzata attraverso una
funzione logica:
\begin{equation}
\label{eq:lead_time}
\tau^* = \min\bigl\{\tau' \in \{1, \ldots, \tau\} : 
  \mathrm{is\_violated}\bigl(\hat{\Phi}_i(t+\tau'),\; \Phi_i^{\text{SLA}}\bigr)\bigr\}
\end{equation}
In questa formula, $\Phi_i^{\text{SLA}}$ rappresenta la soglia stabilita nel modello di
certificazione. A livello operativo, un ipotetico \textit{lead time} di 12 giorni indica che si
dispone di dodici giorni per intervenire sul sistema prima che avvenga la violazione.
Qualora l'insieme risulti vuoto, ovvero nessun passo $\tau'$ nell'orizzonte 
produca una violazione prevista, la Fase~IV non genera alcun \textit{alert} 
per la proprietà $\Phi_i$ nel ciclo corrente.

\subsubsection{Classificazione della Criticità dell'Alert}
\label{sec:classificazione_criticita}

In seguito all'identificazione di una violazione, è necessario determinare la criticità
dell'\textit{alert}. Quest'ultima viene determinata sulla base della combinazione di quattro fattori
distinti: il \textit{lead time} stimato, lo \textit{score} di anomalia prodotto dall'Isolation
Forest nella Fase~III, il segnale di deviazione comportamentale prodotto dal \gls{CUSUM} e il
segnale di deviazione strutturale confermata prodotto dal Validatore strutturale.

\begin{itemize}
  \item \textbf{Livello Giallo (Attenzione):} si attiva quando il \textit{lead time} stimato
        supera i 7 giorni e non si rilevano segnali strutturali attivi da parte del \gls{CUSUM}
        o dell'Isolation Forest. La previsione è coerente con la \textit{baseline} e non vengono
        rilevate anomalie strutturali: la tendenza al degrado richiede monitoraggio, non un
        intervento immediato;
  \item \textbf{Livello Arancione (Prevenzione):} scatta quando il \textit{lead time} è compreso
        tra 2 e 7 giorni, oppure nel caso in cui il \gls{CUSUM} rilevi un \textit{drift}
        strutturale attivo o l'Isolation Forest segnali uno stato aggregato anomalo. In questo
        caso, la tendenza è considerabile come preoccupante e il sistema richiede un intervento
        preventivo;
  \item \textbf{Livello Rosso (Emergenza):} viene generato sempre quando il \textit{lead time}
        stimato è inferiore a 2 giorni. La presenza simultanea della conferma causale di Granger
        sulla catena identificata e del \textit{Path Adherence Score} in rapido declino eleva la
        certezza diagnostica del giudizio. Analogamente, l'attivazione di segnali quali anomalia
        Isolation Forest attiva e il segnale di deviazione strutturale confermata dal Validatore
        strutturale nella Fase~III costituisce un'ulteriore conferma della gravità sistemica, pur
        non essendone condizione necessaria. Tuttavia, in assenza di tali riscontri causali e
        topologici, il livello Rosso scatta ugualmente a protezione del sistema, includendo
        nell'\textit{alert} un indicatore di affidabilità ridotta per segnalare che la violazione
        prevista non è ancora aggravata da derive strutturali manifeste. Questo scenario indica una
        violazione imminente della \textit{compliance}, richiedendo un intervento tecnico immediato.
\end{itemize}

I livelli vengono valutati in ordine decrescente di gravità: il Rosso ha precedenza sull'Arancione,
che a sua volta ha la precedenza sul Giallo. In aggiunta ai criteri sopra descritti, qualora la divergenza tra la
previsione del modello principale e quella della \textit{baseline} lineare superi la soglia configurabile 
definita nella Fase~I, la classificazione dell'\textit{alert} viene ridotta di un livello, da Rosso ad Arancione
oppure da Arancione a Giallo. In tal caso, l'\textit{alert} include anche un indicatore esplicito di affidabilità ridotta, 
segnalando che la stima è soggetta a incertezza modellistica elevata. Il livello Giallo non viene ulteriormente
declassato: la tendenza al degrado richiede comunque monitoraggio indipendentemente dall'incertezza previsionale.

\subsubsection{Struttura dell'Alert Semanticamente Arricchito}
\label{sec:struttura_alert}

L'output finale della Fase~IV, per ogni proprietà a rischio, è un \textit{alert} strutturato che
include una serie di campi semantici, ciascuno derivante dalle analisi compiute nelle fasi
precedenti. I campi che compongono l'\textit{alert} includono: la \textit{proprietà a rischio}, che
identifica la proprietà certificata sotto minaccia; il \textit{compliance set impattato}, che elenca
l'insieme di componenti la cui conformità congiunta è minacciata; il \textit{lead time stimato}
$\tau^*$, che indica il tempo residuo prima che si verifichi la violazione effettiva e la
\textit{metrica critica}, che riporta il valore previsto e la soglia certificata della metrica che
sta guidando la violazione. L'\textit{arco critico principale}, invece, identifica l'arco del grafo
temporale che contribuisce maggiormente alla violazione aggregata della proprietà monitorata. 
La \textit{causa radice} indica il componente o la catena causale identificata dalla Fase~II come origine 
del degrado, mentre l'\textit{interferenza cross-property} nomina la proprietà $\Phi_j$ la cui attività 
contribuisce al degrado di $\Phi_i$, qualora la causa radice coinvolga un nodo condiviso. Infine, 
viene indicato il livello di criticità.

\begin{table}[htbp]
\centering
\caption{Esempio di \textit{alert} strutturato per lo scenario di interferenza \textit{cross-property}
(valori esemplificativi).}
\label{tab:alert_telemed}
\begin{tabular}{ll}
\toprule
\textbf{Campo} & \textbf{Valore (esemplificativo)} \\
\midrule
Proprietà a rischio       & \textit{Real-time Dependability} $\Phi_{\text{crit}}$ \\
\textit{Compliance Set}   & $H_{\text{crit}} = \{\text{SGW, SIN, SDB, SAL}\}$ \\
\textit{Lead time} stimato & $\tau^* = 12$ giorni \\
Metrica critica           & $\hat{L}_{\text{crit}} > 500\;\text{ms}$ (soglia certificata) \\
Arco critico principale   & $e_2 : \text{SIN} \to \text{SDB}$ (contributo prevalente) \\
Causa radice              & Saturazione \textit{connection pool} SDB \\
Interferenza              & $\Phi_{\text{adm}}$: \textit{throughput} $\theta_{e_5}$ in crescita \\
Nodo condiviso            & $\text{SDB} \in \text{Shared}(H_{\text{crit}}, H_{\text{adm}})$ \\
Criticità                 & Arancione (\gls{CUSUM} attivo, \textit{lead time} 12 giorni) \\
\bottomrule
\end{tabular}
\end{table}

\subsubsection{Contributo Distintivo Rispetto agli Approcci Tradizionali}
\label{sec:contributo_distintivo}

Il confronto tra l'\textit{alert} generato dall'approccio ibrido proposto e quello che produrrebbe
un sistema di monitoraggio tradizionale illustra concretamente il valore aggiunto del Layer~1 e
della Fase~IV. Privo di contesto semantico, un sistema puramente statistico si limiterebbe a
segnalare, nello scenario di riferimento, che è previsto un aumento entro un determinato orizzonte
temporale della latenza sull'arco $e_2\!:\!\text{SIN}\!\to\!\text{SDB}$, evidenziando al contempo
un'alta correlazione con la saturazione del \textit{connection pool} di SDB. Seppur tecnicamente
corretto, questo \textit{alert} risulta incompleto sotto tre aspetti cruciali.

Il primo aspetto riguarda la \textit{proprietà certificata}; il sistema non è infatti in grado di
nominare $\Phi_{\text{crit}}$ in quanto non possiede un'associazione esplicita tra l'arco $e_2$ e
la proprietà che esso serve. Sistemi come FIRM~\cite{qiu2020firm}, pur raggiungendo un'accuracy
media del 93\% nella localizzazione della violazione, producono diagnosi che non specificano quale
proprietà certificata sia effettivamente minacciata. Il secondo aspetto riguarda
l'\textit{interferenza cross-property}. Un sistema tradizionale non includerebbe il
\textit{throughput} $\theta_{e_5}$ nella propria analisi perché tale arco appartiene a un flusso
amministrativo distinto, separato a livello logico da quello clinico. Questa limitazione strutturale
viene riscontrata anche in sistemi come GAMMA~\cite{somashekar2024gamma}, progettati per la suite
DeathStarBench~\cite{gan2019deathstar}. In questi casi, il rilevamento di \textit{bottleneck}
multipli avviene a livello di singola richiesta, senza una mappa esplicita delle risorse condivise
tra flussi concorrenti. Il terzo aspetto riguarda la \textit{semantica dell'aggregazione}. Un
sistema tradizionale non calcolerebbe la latenza \textit{end-to-end} prevista in quanto non conosce
la topologia dell'intero percorso e la relativa funzione di aggregazione. Di conseguenza,
l'eventuale superamento della soglia di 500\,ms risulterebbe evidente solo nel momento in cui
l'operatore aggregasse manualmente le latenze dei tre archi del percorso clinico.

L'approccio ibrido proposto risolve queste tre limitazioni. Il \textit{mapping} semantico
(Layer~3) fornisce il contesto associando $e_2$ a $H_{\text{crit}}$, la funzione logica
dell'ipergrafo intercetta le cause esterne, includendo così $\theta_{e_5}$ nell'analisi
dell'interferenza, e la Fase~IV applica proattivamente la funzione di aggregazione additiva. Il
valore aggiunto non è solo operativo ma anche computazionale; la selezione topologica delle metriche
in Fase~I e lo \textit{screening} topologico delle coppie in Fase~II riducono lo spazio di analisi
da quadratico nel numero totale di metriche del sistema a lineare nel prodotto tra il numero di
proprietà monitorate e la dimensione media dei \textit{compliance sets}. Entrambe le riduzioni sono
conseguenze dirette dell'adozione dell'ipergrafo, per questo motivo sono garantite
indipendentemente dalla qualità e dalla quantità dei dati storici a disposizione.